% Generated by the blueprint tool (blueprint-pdf/1) from a blueprint snapshot; edit the
% Markdown sources, not this file.  Compile with `latexmk -pdf` or by running
% `pdflatex` twice (the second run settles the contents and the page references).
% Images, if any, are read from assets/ next to this file.
\documentclass[11pt]{article}
\usepackage[T1]{fontenc}
\usepackage[utf8]{inputenc}
\usepackage{lmodern}
\usepackage{textcomp}
\usepackage[a4paper,margin=2.5cm,headheight=14pt]{geometry}
\usepackage{amsmath,amssymb,amsthm}
\usepackage{amscd}  % KaTeX's CD diagrams; part of amsmath, so always there
\usepackage{mathtools}
\usepackage{bm}
\usepackage{graphicx}
\usepackage{xcolor}
\usepackage{array,booktabs,tabularx}
\usepackage{fancyvrb}
\usepackage[normalem]{ulem}
\usepackage{microtype}
% Optional, and only with their fonts: a small TeX installation has the
% packages without the fonts, which fails at the first formula.
\IfFileExists{mathrsfs.sty}{\IfFileExists{rsfs10.tfm}{\usepackage{mathrsfs}}{}}{}
\IfFileExists{stmaryrd.sty}{\IfFileExists{stmary10.tfm}{\usepackage{stmaryrd}}{}}{}
\IfFileExists{cancel.sty}{\usepackage{cancel}}{}
\usepackage[hyphens]{url}
\usepackage[pdfusetitle,bookmarksnumbered=false]{hyperref}
\usepackage{bookmark}
\definecolor{bplink}{HTML}{1F4E9C}
\definecolor{bpmuted}{HTML}{5F676E}
\hypersetup{colorlinks=true,linkcolor=bplink,urlcolor=bplink,citecolor=bplink}
\setlength{\parindent}{0pt}
\setlength{\parskip}{0.55em plus 0.2em minus 0.1em}
% Lists as tight as the page's.
\makeatletter
\def\@listi{\leftmargin\leftmargini\parsep 0pt\topsep 0.3em plus 0.1em\itemsep 0.1em plus 0.1em}
\let\@listI\@listi
\def\@listii{\leftmargin\leftmarginii\labelwidth\leftmarginii\advance\labelwidth-\labelsep\parsep 0pt\topsep 0.1em\itemsep 0.1em}
% An ordered list that does not start at 1.
\newcommand{\bpstartat}[1]{\setcounter{\@enumctr}{\numexpr#1-1\relax}}
\makeatother
\setcounter{tocdepth}{3}
\emergencystretch=3em
\allowdisplaybreaks

% A numbered statement, set as a paper sets one: "Theorem 2.1 (Title)."
% #1 is the lead ("Theorem~2.1"), #2 what follows it in roman.
\newenvironment{bpstatement}[2]{\par\addvspace{0.6\baselineskip}\noindent\textbf{#1}#2\textbf{.}\enspace\ignorespaces}
  {\par\addvspace{0.4\baselineskip}}
\newenvironment{bpstatementit}[2]{\begin{bpstatement}{#1}{#2}\itshape\ignorespaces}{\end{bpstatement}}
% A step: a relation with prose of its own, read under the object it starts from.
\newenvironment{bpstep}[1]{\par\addvspace{0.3\baselineskip}\begin{list}{}{\leftmargin=1.6em\rightmargin=0pt\topsep=0pt}\item\relax
  {\small\color{bpmuted}#1}\par\nopagebreak}{\end{list}\par}
% Notes about where an object is written out.
\newenvironment{bpnote}{\par\small\color{bpmuted}\itshape}{\par}
\newcommand{\bpomitted}[1]{\par{\small\color{bpmuted}\itshape #1\par}}
\newcommand{\bpbodyheading}[1]{\par\addvspace{0.4\baselineskip}\noindent\textbf{#1}\par\nopagebreak}
\newcommand{\bprule}{\par\medskip\noindent\hfil\rule{0.3\linewidth}{0.4pt}\hfil\par\medskip}
\newcommand{\bpboundary}[1]{\par\nopagebreak{\small\color{bpmuted}#1\par}\nopagebreak}
% A reference to something this export leaves out, and one that resolves nowhere.
\newcommand{\bpout}[1]{#1\textsuperscript{\hyperlink{bp-outside}{\ensuremath{\ast}}}}
\newcommand{\bpbroken}[1]{\textcolor{red!70!black}{[#1?]}}
\newcommand{\bpslug}[1]{\texttt{#1}}
% A bibliography entry: label, anchor, entry; the label hangs in the margin.
\newcommand{\bpbibitem}[3]{\par\setbox0\hbox{[#1]\enspace}\ifdim\wd0<5em\setbox0\hbox to 5em{[#1]\hfil}\fi
  \hangindent=\wd0 \hangafter=1 \noindent\box0\phantomsection\label{#2}#3\par\smallskip}
% A character the export knows no way to set.
\newcommand{\bpglyph}[1]{\fbox{\scriptsize #1}}
% Images, at their own size unless they are wider than the line.
\newsavebox{\bpimagebox}
\newcommand{\bpimage}[1]{\sbox{\bpimagebox}{\includegraphics{#1}}%
  \ifdim\wd\bpimagebox>\linewidth\resizebox{\linewidth}{!}{\usebox{\bpimagebox}}\else\usebox{\bpimagebox}\fi}
\newcommand{\bpmissingimage}[1]{\fbox{\parbox{0.8\linewidth}{\centering\small\color{bpmuted}Image not included: #1}}}
\DefineVerbatimEnvironment{bpverbatim}{Verbatim}{fontsize=\small,xleftmargin=1em}
\renewcommand{\qedsymbol}{\ensuremath{\blacksquare}}
\hypersetup{pdftitle={Weil Conjectures — 3 Constructible complexes},pdfsubject={Chapter 3 Constructible complexes; all levels},pdfcreator={blueprint}}
% KaTeX commands the website knows
\providecommand{\R}{\mathbb{R}}\providecommand{\N}{\mathbb{N}}\providecommand{\Z}{\mathbb{Z}}
\providecommand{\Q}{\mathbb{Q}}\providecommand{\C}{\mathbb{C}}
\providecommand{\Reals}{\mathbb{R}}\providecommand{\reals}{\mathbb{R}}\providecommand{\Complex}{\mathbb{C}}
\providecommand{\cnums}{\mathbb{C}}\providecommand{\natnums}{\mathbb{N}}\providecommand{\Bbb}[1]{\mathbb{#1}}
\providecommand{\bold}[1]{\mathbf{#1}}\providecommand{\frak}[1]{\mathfrak{#1}}\providecommand{\mathbbm}[1]{\mathbb{#1}}
\providecommand{\lang}{\langle}\providecommand{\rang}{\rangle}\providecommand{\bra}[1]{\langle #1\rvert}
\providecommand{\ket}[1]{\lvert #1\rangle}\providecommand{\braket}[1]{\langle #1\rangle}
\providecommand{\set}[1]{\{#1\}}\providecommand{\Set}[1]{\left\{#1\right\}}
\providecommand{\infin}{\infty}\providecommand{\isin}{\in}\providecommand{\plusmn}{\pm}\providecommand{\sub}{\subset}
\providecommand{\sube}{\subseteq}\providecommand{\supe}{\supseteq}\providecommand{\exist}{\exists}
\providecommand{\bull}{\bullet}\providecommand{\larr}{\leftarrow}\providecommand{\rarr}{\rightarrow}
\providecommand{\lrarr}{\leftrightarrow}\providecommand{\Larr}{\Leftarrow}\providecommand{\Rarr}{\Rightarrow}
\providecommand{\Lrarr}{\Leftrightarrow}\providecommand{\uarr}{\uparrow}\providecommand{\darr}{\downarrow}
\providecommand{\harr}{\leftrightarrow}\providecommand{\lArr}{\Leftarrow}\providecommand{\rArr}{\Rightarrow}
\providecommand{\hArr}{\Leftrightarrow}\providecommand{\alef}{\aleph}\providecommand{\alefsym}{\aleph}
\providecommand{\real}{\Re}\providecommand{\image}{\Im}\providecommand{\weierp}{\wp}
\providecommand{\Alpha}{\mathrm{A}}\providecommand{\Beta}{\mathrm{B}}\providecommand{\Epsilon}{\mathrm{E}}
\providecommand{\Zeta}{\mathrm{Z}}\providecommand{\Eta}{\mathrm{H}}\providecommand{\Iota}{\mathrm{I}}
\providecommand{\Kappa}{\mathrm{K}}\providecommand{\Mu}{\mathrm{M}}\providecommand{\Nu}{\mathrm{N}}
\providecommand{\Omicron}{\mathrm{O}}\providecommand{\Rho}{\mathrm{P}}\providecommand{\Tau}{\mathrm{T}}
\providecommand{\Chi}{\mathrm{X}}\providecommand{\omicron}{o}\providecommand{\gt}{>}\providecommand{\lt}{<}
\providecommand{\And}{\&}\providecommand{\hdashline}{\hline}\providecommand{\operatornamewithlimits}[1]{\operatorname*{#1}}
\providecommand{\llbracket}{[\![}\providecommand{\rrbracket}{]\!]}\providecommand{\cancel}[1]{#1}
\providecommand{\mathscr}{\mathcal}
% the characters beyond ASCII this document uses
\DeclareUnicodeCharacter{2160}{\textup{I}}
\DeclareUnicodeCharacter{2161}{\textup{II}}
\DeclareUnicodeCharacter{2162}{\textup{III}}
% the project's macros ([katex.macros] in blueprint.toml)
\providecommand{\A}{}\renewcommand{\A}{\mathbf A}
\providecommand{\AEn}{}\renewcommand{\AEn}{A_{E,n}}
\providecommand{\Aaf}{}\renewcommand{\Aaf}{\mathbf A}
\providecommand{\Ab}{}\renewcommand{\Ab}{\mathsf{Ab}}
\providecommand{\AffSch}{}\renewcommand{\AffSch}{\operatorname{AffSch}}
\providecommand{\Aone}{}\renewcommand{\Aone}{\mathbf A^1}
\providecommand{\Bl}{}\renewcommand{\Bl}{\operatorname{Bl}}
\providecommand{\Br}{}\renewcommand{\Br}{\mathrm{Br}}
\providecommand{\C}{}\renewcommand{\C}{\mathbb C}
\providecommand{\Cech}{}\renewcommand{\Cech}{\check{H}}
\providecommand{\Coext}{}\renewcommand{\Coext}{\operatorname{Coext}}
\providecommand{\Cone}{}\renewcommand{\Cone}{\operatorname{Cone}}
\providecommand{\Cons}{}\renewcommand{\Cons}{\operatorname{Cons}}
\providecommand{\Cref}{}\renewcommand{\Cref}[1]{\text{#1}}
\providecommand{\D}{}\renewcommand{\D}{D}
\providecommand{\DLAFT}{}\renewcommand{\DLAFT}{\operatorname{FT}}
\providecommand{\DLAFr}{}\renewcommand{\DLAFr}{\operatorname{Fr}}
\providecommand{\DLAQ}{}\renewcommand{\DLAQ}{{\overline{\mathbf Q}_\ell}}
\providecommand{\DLRQ}{}\renewcommand{\DLRQ}{{\overline{\mathbf Q}_\ell}}
\providecommand{\Dbc}{}\renewcommand{\Dbc}{D^b_c}
\providecommand{\Dc}{}\renewcommand{\Dc}{D_{\mathrm{cons}}}
\providecommand{\Dctf}{}\renewcommand{\Dctf}{D^b_{ctf}}
\providecommand{\Dic}{}\renewcommand{\Dic}{D_{\IndCons}}
\providecommand{\Dplus}{}\renewcommand{\Dplus}{D^{+}}
\providecommand{\End}{}\renewcommand{\End}{\operatorname{End}}
\providecommand{\Et}{}\renewcommand{\Et}{\mathrm{\acute{e}t}}
\providecommand{\Ex}{}\renewcommand{\Ex}{\mathrm{Ex}}
\providecommand{\Ext}{}\renewcommand{\Ext}{\operatorname{Ext}}
\providecommand{\F}{}\renewcommand{\F}{\mathbb F}
\providecommand{\FT}{}\renewcommand{\FT}{\operatorname{FT}}
\providecommand{\Fbar}{}\renewcommand{\Fbar}{{\overline{\mathbf F}_q}}
\providecommand{\Fell}{}\renewcommand{\Fell}{\mathbf F_\ell}
\providecommand{\Fib}{}\renewcommand{\Fib}{\operatorname{Fib}}
\providecommand{\Fix}{}\renewcommand{\Fix}{\operatorname{Fix}}
\providecommand{\Fp}{}\renewcommand{\Fp}{{\mathbf F_{p}}}
\providecommand{\Fq}{}\renewcommand{\Fq}{{\mathbf F_q}}
\providecommand{\Fqn}{}\renewcommand{\Fqn}{{\mathbf F_{q^{n}}}}
\providecommand{\Frob}{}\renewcommand{\Frob}{\operatorname{Frob}}
\providecommand{\GL}{}\renewcommand{\GL}{\operatorname{GL}}
\providecommand{\Gal}{}\renewcommand{\Gal}{\operatorname{Gal}}
\providecommand{\Hens}{}\renewcommand{\Hens}{\operatorname{Hens}}
\providecommand{\Hilb}{}\renewcommand{\Hilb}{\operatorname{Hilb}}
\providecommand{\Hom}{}\renewcommand{\Hom}{\operatorname{Hom}}
\providecommand{\IndCons}{}\renewcommand{\IndCons}{\operatorname{IndCons}}
\providecommand{\Kar}{}\renewcommand{\Kar}{\operatorname{Kar}}
\providecommand{\LT}{}\renewcommand{\LT}{\operatorname{LT}}
\providecommand{\Lie}{}\renewcommand{\Lie}{\operatorname{Lie}}
\providecommand{\Loc}{}\renewcommand{\Loc}{\operatorname{Loc}}
\providecommand{\Mod}{}\renewcommand{\Mod}{\mathsf{Mod}}
\providecommand{\N}{}\renewcommand{\N}{\mathbb{N}}
\providecommand{\Needspace}{}\renewcommand{\Needspace}{\par}
\providecommand{\O}{}\renewcommand{\O}{{\mathcal{O}}}
\providecommand{\Ob}{}\renewcommand{\Ob}{\mathrm{Ob}}
\providecommand{\Oe}{}\renewcommand{\Oe}{\mathcal O_E}
\providecommand{\PP}{}\renewcommand{\PP}{\mathbf{P}}
\providecommand{\Perf}{}\renewcommand{\Perf}{\operatorname{Perf}}
\providecommand{\Pic}{}\renewcommand{\Pic}{\mathrm{Pic}}
\providecommand{\Pm}{}\renewcommand{\Pm}{\mathbf P^{m}}
\providecommand{\Pone}{}\renewcommand{\Pone}{\mathbf P^{1}}
\providecommand{\Pqn}{}\renewcommand{\Pqn}{\mathbf P^{n}}
\providecommand{\ProFin}{}\renewcommand{\ProFin}{\operatorname{ProFin}}
\providecommand{\Proet}{}\renewcommand{\Proet}{\mathrm{pro\acute{e}t}}
\providecommand{\Q}{}\renewcommand{\Q}{\mathbb Q}
\providecommand{\Qb}{}\renewcommand{\Qb}{\overline{\mathbf Q}_{\ell}}
\providecommand{\Qbar}{}\renewcommand{\Qbar}{\overline{\mathbf Q}}
\providecommand{\QbarX}{}\renewcommand{\QbarX}{\overline{\mathbf Q}_{\ell,X}}
\providecommand{\Qbarl}{}\renewcommand{\Qbarl}{{\overline{\mathbf Q}_{\ell}}}
\providecommand{\Qell}{}\renewcommand{\Qell}{{\mathbf{Q}_{\ell}}}
\providecommand{\Ql}{}\renewcommand{\Ql}{\mathbf Q_{\ell}}
\providecommand{\Qlb}{}\renewcommand{\Qlb}{{\overline{\mathbf Q}_\ell}}
\providecommand{\Qlbar}{}\renewcommand{\Qlbar}{\overline{\mathbf Q}_\ell}
\providecommand{\R}{}\renewcommand{\R}{\mathbb{R}}
\providecommand{\RG}{}\renewcommand{\RG}{\operatorname{\mathrm{R}\Gamma}}
\providecommand{\RGamma}{}\renewcommand{\RGamma}{R\Gamma}
\providecommand{\RGammac}{}\renewcommand{\RGammac}{R\Gamma_c}
\providecommand{\RHom}{}\renewcommand{\RHom}{R\mathcal{H}om}
\providecommand{\RHomX}{}\renewcommand{\RHomX}{R\mathcal H om_X}
\providecommand{\Rat}{}\renewcommand{\Rat}{\operatorname{IsRat}}
\providecommand{\Res}{}\renewcommand{\Res}{\operatorname{Res}}
\providecommand{\Rlim}{}\renewcommand{\Rlim}{R\!\varprojlim}
\providecommand{\SIm}{}\renewcommand{\SIm}{\operatorname{SIm}}
\providecommand{\Sch}{}\renewcommand{\Sch}{\mathsf{Sch}}
\providecommand{\Sh}{}\renewcommand{\Sh}{\operatorname{Sh}}
\providecommand{\Span}{}\renewcommand{\Span}{\operatorname{span}}
\providecommand{\Spec}{}\renewcommand{\Spec}{\operatorname{Spec}}
\providecommand{\Tr}{}\renewcommand{\Tr}{\operatorname{Tr}}
\providecommand{\Trsm}{}\renewcommand{\Trsm}{\operatorname{Tr}^{\mathrm{sm}}}
\providecommand{\Wt}{}\renewcommand{\Wt}{\operatorname{Wt}}
\providecommand{\Z}{}\renewcommand{\Z}{\mathbb{Z}}
\providecommand{\Zar}{}\renewcommand{\Zar}{\mathrm{zar}}
\providecommand{\Zell}{}\renewcommand{\Zell}{{\mathbf{Z}_{\ell}}}
\providecommand{\Zl}{}\renewcommand{\Zl}{{\mathbf Z_\ell}}
\providecommand{\affet}{}\renewcommand{\affet}[1]{{#1}^{\mathrm{aff}}_{\mathrm{ét}}}
\providecommand{\affproet}{}\renewcommand{\affproet}[1]{{#1}^{\mathrm{aff}}_{\mathrm{proét}}}
\providecommand{\boxt}{}\renewcommand{\boxt}{\boxtimes_A^{\mathrm L}}
\providecommand{\cF}{}\renewcommand{\cF}{\mathcal{F}}
\providecommand{\cHom}{}\renewcommand{\cHom}{R\mathcal{H}\!om}
\providecommand{\cI}{}\renewcommand{\cI}{\mathcal{I}}
\providecommand{\cL}{}\renewcommand{\cL}{\mathcal L}
\providecommand{\cM}{}\renewcommand{\cM}{\mathcal M}
\providecommand{\cN}{}\renewcommand{\cN}{\mathcal N}
\providecommand{\cO}{}\renewcommand{\cO}{\mathcal O}
\providecommand{\cT}{}\renewcommand{\cT}{\mathcal T}
\providecommand{\checked}{}\renewcommand{\checked}{\color{red}\aftergroup\restorecolor}
\providecommand{\coeff}{}\renewcommand{\coeff}{\operatorname{coeff}}
\providecommand{\coker}{}\renewcommand{\coker}{\operatorname{coker}}
\providecommand{\colim}{}\renewcommand{\colim}{\operatorname{colim}}
\providecommand{\constructedby}{}\renewcommand{\constructedby}[1]{\par\smallskip\noindent\textcolor{purple!60!black}{\textsf{Constructed by:} }
  \footnotesize\sloppy\nolinkurl{#1}\par}
\providecommand{\deepresult}{}\renewcommand{\deepresult}[1]{\par\smallskip\noindent\textcolor{orange!45!black}{
  \textsf{Blueprint role: named deep theorem node; see the cluster proof architecture.}}
  \par\noindent\textcolor{orange!45!black}{\textsf{Mathematical source:} }#1\par}
\providecommand{\dt}{}\renewcommand{\dt}{\mathbin{\otimes^{\mathbf L}}}
\providecommand{\dw}{}\renewcommand{\dw}{\operatorname{dw}}
\providecommand{\et}{}\renewcommand{\et}[1]{{#1}_{\mathrm{ét}}}
\providecommand{\ev}{}\renewcommand{\ev}{\operatorname{ev}}
\providecommand{\fib}{}\renewcommand{\fib}{\operatorname{fib}}
\providecommand{\formalizationnote}{}\renewcommand{\formalizationnote}[1]{\begin{quote}\small\raggedright\textbf{Formalization note.} #1\end{quote}}
\providecommand{\fpqc}{}\renewcommand{\fpqc}{X_{\mathrm{fpqc}}}
\providecommand{\gammacont}{}\renewcommand{\gammacont}{\Gamma_{\mathrm{cont}}}
\providecommand{\gm}{}\renewcommand{\gm}{\mathbf G_m}
\providecommand{\hocolim}{}\renewcommand{\hocolim}{\operatorname{hocolim}}
\providecommand{\holim}{}\renewcommand{\holim}{\operatorname{holim}}
\providecommand{\hrB}{}\renewcommand{\hrB}{\@ifnextchar]{\hrEnd}{\hrStep}}
\providecommand{\hrEnd}{}\renewcommand{\hrEnd}{\@firstoftwo{}}
\providecommand{\hrStep}{}\renewcommand{\hrStep}{\@firstoftwo{\hrB}}
\providecommand{\id}{}\renewcommand{\id}{\operatorname{id}}
\providecommand{\kunFl}{}\renewcommand{\kunFl}{\mathbf F_\ell}
\providecommand{\kunZn}{}\renewcommand{\kunZn}{\mathbf Z/n\mathbf Z}
\providecommand{\kuncd}{}\renewcommand{\kuncd}{\operatorname{cd}}
\providecommand{\kundim}{}\renewcommand{\kundim}{\dim}
\providecommand{\kunotimes}{}\renewcommand{\kunotimes}[1]{\otimes^{\mathbf L}_{#1}}
\providecommand{\kunsupp}{}\renewcommand{\kunsupp}{\operatorname{Supp}}
\providecommand{\lotimes}{}\renewcommand{\lotimes}{\otimes^{\mathrm L}}
\providecommand{\m}{}\renewcommand{\m}{{\mathfrak{m}}}
\providecommand{\n}{}\renewcommand{\n}{{\mathfrak{n}}}
\providecommand{\nagataGap}{}\renewcommand{\nagataGap}[1]{\par\noindent\textbf{Formalization gap.} #1\par}
\providecommand{\openobligation}{}\renewcommand{\openobligation}[1]{\par\smallskip\noindent\textcolor{red!65!black}{\textsf{Open obligation:} }
  \footnotesize\sloppy\nolinkurl{#1}\par}
\providecommand{\p}{}\renewcommand{\p}{{\mathfrak{p}}}
\providecommand{\pr}{}\renewcommand{\pr}{\operatorname{pr}}
\providecommand{\preshv}{}\renewcommand{\preshv}[1]{\mathrm{PreShv}(#1)}
\providecommand{\pro}{}\renewcommand{\pro}[1]{\mathrm{Pro}(#1)}
\providecommand{\proet}{}\renewcommand{\proet}{X_{\mathrm{proét}}}
\providecommand{\q}{}\renewcommand{\q}{{\mathfrak{q}}}
\providecommand{\rad}{}\renewcommand{\rad}{\operatorname{rad}}
\providecommand{\rank}{}\renewcommand{\rank}{\operatorname{rank}}
\providecommand{\red}{}\renewcommand{\red}{\color{red}}
\providecommand{\restorecolor}{}\renewcommand{\restorecolor}{\color{black}}
\providecommand{\retiredobligation}{}\renewcommand{\retiredobligation}[1]{\par\smallskip\noindent\textcolor{red!75!black}{\textsf{Retired invalid route:} }
  \footnotesize\sloppy\nolinkurl{#1}\par}
\providecommand{\rk}{}\renewcommand{\rk}{\operatorname{rk}}
\providecommand{\sep}{}\renewcommand{\sep}{\textnormal{sep}}
\providecommand{\shv}{}\renewcommand{\shv}[1]{\mathrm{Shv}(#1)}
\providecommand{\spec}{}\renewcommand{\spec}[1]{\mathrm{Spec}(#1)}
\providecommand{\stacks}{}\renewcommand{\stacks}[1]{\href{https://stacks.math.columbia.edu/tag/#1}{\textnormal{Stacks, Tag~#1}}}
\providecommand{\stacksproject}{}\renewcommand{\stacksproject}[1]{{\cite[\href{https://stacks.math.columbia.edu/tag/#1}{Tag #1}]{stacks-project}}}
\providecommand{\statusnote}{}\renewcommand{\statusnote}[1]{\begin{quote}\small\raggedright\textbf{Status.} #1\end{quote}}
\providecommand{\tildeideal}{}\renewcommand{\tildeideal}[2]{\widetilde{(#1, #2)}}
\providecommand{\tildering}{}\renewcommand{\tildering}[3]{{#1}_{\widetilde{(#2, #3)}}}
\providecommand{\tr}{}\renewcommand{\tr}{\operatorname{tr}}
\providecommand{\wt}{}\renewcommand{\wt}{\operatorname{wt}}

\begin{document}
\pagestyle{plain}
\begin{center}
{\LARGE\bfseries Weil Conjectures\par}
\vspace{0.7em}
{\large Chapter 3: Constructible complexes\par}
\vspace{0.5em}
{\small\color{bpmuted}Snapshot 2026-09-26, revision 8e4cb88e2a\par}
\end{center}
\vspace{0.5em}
{\small\tableofcontents}
\bigskip\hrule\bigskip

% object: chap-Constructible-complexes (blueprint/chap-Constructible-complexes/_section.md)
\section*{\texorpdfstring{3\quad Constructible complexes}{3 Constructible complexes}}
\phantomsection\label{bp1}%
\addcontentsline{toc}{section}{\texorpdfstring{3\quad Constructible complexes}{3 Constructible complexes}}
% text: prose-03-01
\phantomsection\label{bp2}%
A lisse \(E\)-sheaf is the analogue of a local system: pro-étale locally it is isomorphic to \(E_X^r\) for some \(r\) (Definition \hyperref[bp9]{3.2} below). Inside such a sheaf \(L\), a lattice is a locally free \(\mathcal O_E\)-subsheaf \(L_0\) with \(L_0[1/\ell]=L\), and the first point of this chapter is that a lattice need not exist. For an example, glue \(0\) and \(\infty\) on \(\mathbf P^1\) to obtain a nodal curve \(X\). Then the infinite chain \(Y\) of projective lines, with \(\infty\) on each line glued to \(0\) on the next, maps to \(X\) by an étale covering. Moreover, the group \(\mathbf Z\) acts on \(Y\) by translation, simply transitively on the fibres, so \(Y\to X\) is an étale \(\mathbf Z\)-torsor. Using this torsor and letting the generator of \(\mathbf Z\) act on \(E\) by multiplication by \(\ell\), we obtain a lisse \(E\)-sheaf of rank one. It has no lattice, because \(\ell\) is not a unit of \(\mathcal O_E\), so that multiplication by \(\ell\) does not carry any lattice onto itself. The sheaf comes from the character \(n\mapsto\ell^n\) of \(\mathbf Z\), and since this character does not extend continuously to \(\widehat{\mathbf Z}\), the sheaf has no classical counterpart as a system \((F_n)\) (see the discussion before Lemma \hyperref[bp77]{3.24}; the example is in [\hyperref[bp93]{BS15}]).

This matters because every geometric input to the next chapters (proper base change, the cohomology of curves, Deligne's generic base change) is a theorem about torsion sheaves, whereas a rational sheaf has no reduction modulo \(\ell\): since \(E_X\) is \(\ell\)-divisible, we have \(E_X\otimes^{\mathbf L}\mathcal O_E/\pi_E=0\). To use the torsion theory we must therefore find integral objects inside rational ones, and then pass to their torsion reductions. The example above shows that this cannot be done globally.

Classically each coefficient ring comes with its own theory:

\begin{center}\small
\begin{tabularx}{\linewidth}{>{\raggedright\arraybackslash}X>{\raggedright\arraybackslash}X}
\toprule
coefficients & classical constructible objects \\
\midrule
finite ring \(A\) & constructible étale sheaves of \(A\)-modules \\
\(\mathbf Z_\ell\), \(\mathcal O_E\) & systems \((F_n)\) of constructible   \(\mathcal O_E/\pi_E^n\)-sheaves \\
 & with \(F_{n+1}\otimes\mathcal O_E/\pi_E^n\simeq F_n\) \\
\(\mathbf Q_\ell\), \(E\) & the same systems up to isogeny, that is, with \(\ell\) inverted \\
\(\overline{\mathbf Q}_\ell\) & a colimit of the categories for all finite \(E/\mathbf Q_\ell\) \\
\bottomrule
\end{tabularx}
\end{center}

On the pro-étale site, by contrast, one local condition covers all of them. Namely, a \(\Lambda_X\)-module is constructible if, on the strata of a finite stratification of \(X\), it is pro-étale locally of the form \(\Lambda_X\otimes_{\Lambda(*)}P\) for a finitely presented module \(P\) over \(\Lambda(*)\), the ring of sections over the point (Definition \hyperref[bp7]{3.1}). For finite \(A\) these are the classical constructible sheaves, while for \(E\) and \(\overline{\mathbf Q}_\ell\) they are lisse on the strata. For \(\mathcal O_E\), on the other hand, the modules \(P\) may have torsion, for example \(P=\mathcal O_E/\pi_E\).

The six operations will act on \(D^b_{\mathrm{ctf}}(X,\Lambda)\), bounded complexes with constructible cohomology and finite Tor amplitude. For \(\Lambda\) not a finite ring the Tor condition is automatic, that is, \(D^b_{\mathrm{ctf}}(X,\Lambda)=\Dbc(X,\Lambda)\) (Proposition \hyperref[bp40]{3.12}). For a finite ring, however, it matters: over \(\mathbf Z/\ell^2\) the constructible sheaf \(\mathbf F_\ell\) has infinite Tor dimension (Section \hyperref[bp5]{3.2} below).

Over \(\mathcal O_E\) the passage to torsion is a criterion: for Noetherian \(X\), a complex is constructible exactly when it is derived complete and its reduction modulo \(\pi_E\) is constructible (Proposition \hyperref[bp34]{3.11}). The criterion is not formal, because the local models of the reductions modulo \(\pi_E^n\) must be made compatible in \(n\) (Proposition \hyperref[bp31]{3.10}, Section \hyperref[bp13]{3.3}). This description is due to Bhatt and Scholze [\hyperref[bp93]{BS15}].

Over \(E\) and \(\overline{\mathbf Q}_\ell\), as the nodal curve shows, lattices exist only after stratifying, and this is enough. After descending from \(\overline{\mathbf Q}_\ell\) to a finite extension \(E\) of \(\mathbf Q_\ell\) and refining a stratification, a lisse sheaf has a lattice on each stratum (Lemma \hyperref[bp77]{3.24}), and the reductions of the lattice modulo \(\pi_E^n\) are classical torsion sheaves (Lemma \hyperref[bp82]{3.25}).

On a connected normal scheme no stratification is needed: there a lattice exists globally, and lisse \(E\)-sheaves are continuous representations of \(\pi_1(X,\bar x)\) (Proposition \bpout{13.10}, proved in Part III).

Dévissage along strata also shows that \(\Hom\) out of a constructible complex commutes with filtered colimits of complexes bounded below by a common degree (Lemma \hyperref[bp68]{3.22}); this continuity is used, together with descent to a finite \(E\), in Chapter \bpout{5}.

Throughout this chapter schemes are again separated and of finite type over \(k\), as in our conventions.

% object: sec-Derived-functors (blueprint/chap-Constructible-complexes/sec-Derived-functors/_section.md), content hidden
\subsection*{\texorpdfstring{3.1\quad Derived functors}{3.1 Derived functors}}
\phantomsection\label{bp3}%
\addcontentsline{toc}{subsection}{\texorpdfstring{3.1\quad Derived functors}{3.1 Derived functors}}
% text: prose-03-02
\phantomsection\label{bp4}%
Four of the six operations, \(f^*\), \(Rf_*\), \(\otimes^{\mathbf L}\) and internal Hom, exist for every condensed ring \(\Lambda\) on the whole unbounded derived categories by the usual resolutions. The remaining two need more work: \(Rf_!\) is built from proper direct image and extension by zero (Section \bpout{5.4}), and its right adjoint comes from Brown representability.

Sheaves of \(\Lambda_X\)-modules can be viewed as module objects over the ring object \(\Lambda_X\) in \(\Sh(X_{\mathrm{pro}\acute{e}t},\operatorname{Ab})\), with the tensor product over \(\Lambda_X\) given by the coequalizer of the two actions \(M\otimes\Lambda_X\otimes N\rightrightarrows M\otimes N\). The category of \(\Lambda_X\)-modules is Grothendieck abelian. Hence complexes have K-injective resolutions, which define \(Rf_*\) and \(\cHom\) [Stacks, \href{https://stacks.math.columbia.edu/tag/07A5}{07A5}, \href{https://stacks.math.columbia.edu/tag/08J7}{08J7}]. They also have K-flat resolutions, which define \(\otimes^{\mathbf L}\), and \(K\otimes^{\mathbf L}-\) is left adjoint to \(\cHom(K,-)\) [Stacks, \href{https://stacks.math.columbia.edu/tag/06YL}{06YL}].

The pullback of a \(\Lambda_Y\)-module \(M\) along \(f:X\to Y\) is \(f^*M=\Lambda_X\otimes_{f^{-1}\Lambda_Y}f^{-1}M\). Since \(f^{-1}\Lambda_Y\simeq\Lambda_X\) (Lemma \bpout{2.54}), this is just \(f^{-1}M\). Therefore \(f^*:\operatorname{Mod}(Y_{\mathrm{pro}\acute{e}t},\Lambda)\to
\operatorname{Mod}(X_{\mathrm{pro}\acute{e}t},\Lambda)\) is exact, and on derived categories it is left adjoint to \(Rf_*:D(X_{\mathrm{pro}\acute{e}t},\Lambda)\to
D(Y_{\mathrm{pro}\acute{e}t},\Lambda)\). For an open immersion \(j:U\to X\), extension by zero \(j_!:\operatorname{Mod}(U_{\mathrm{pro}\acute{e}t},\Lambda)\to\operatorname{Mod}(X_{\mathrm{pro}\acute{e}t},\Lambda)\) is exact and left adjoint to \(j^*\).

Constructible sheaves will be built from lisse sheaves on the strata of a stratification, each extended by zero to all of \(X\) (Lemma \hyperref[bp56]{3.18} and Corollary \hyperref[bp61]{3.20}). We fix this extension by zero now. Let \(a:S\hookrightarrow X\) be a locally closed immersion, for instance the inclusion of a stratum. Then \(S\) is open in its closure, so \(a\) factors as \(a=i\circ j\), with \(j:S\hookrightarrow\overline S\) the open immersion into the closure and \(i:\overline S\hookrightarrow X\) the closed immersion. We put

\[
a_!=i_*j_!:\operatorname{Mod}(S_{\mathrm{pro}\acute{e}t},\Lambda)\longrightarrow\operatorname{Mod}(X_{\mathrm{pro}\acute{e}t},\Lambda).
\]

For a \(\Lambda_S\)-module \(L\), the sheaf \(a_!L\) restricts to \(L\) on \(S\), and its restrictions to \(\overline S\setminus S\) and to \(X\setminus\overline S\) are zero. Indeed \(a^*a_!=j^*i^*i_*j_!=\mathrm{id}\), the sheaf \(j_!L\) vanishes on \(\overline S\setminus S\), and \(i_*\) of a sheaf vanishes on \(X\setminus\overline S\) (Proposition \bpout{2.58}, which applies because the open immersions involved are quasi-compact, \(X\) being Noetherian). The same proposition shows that \(j_!\) and \(i_*\) are exact, so \(a_!\) is exact. It therefore acts on complexes term by term and defines \(a_!:D(S_{\mathrm{pro}\acute{e}t},\Lambda)\to D(X_{\mathrm{pro}\acute{e}t},\Lambda)\) with no derived functor needed. Once \(Rf_!\) is constructed (Section \bpout{5.4}), this functor is \(Ra_!\).


% object: sec-Constructibility (blueprint/chap-Constructible-complexes/sec-Constructibility/_section.md), content hidden
\subsection*{\texorpdfstring{3.2\quad Constructibility}{3.2 Constructibility}}
\phantomsection\label{bp5}%
\addcontentsline{toc}{subsection}{\texorpdfstring{3.2\quad Constructibility}{3.2 Constructibility}}
% text: prose-03-03
\phantomsection\label{bp6}%
With finite coefficients a constructible sheaf is, on the strata of a finite stratification, locally constant with finite stalks. The six operations, however, will act on constructible complexes of finite Tor amplitude for every coefficient ring, so we need a definition of constructibility that covers all coefficient rings at once. We give it by a local condition (Definition \hyperref[bp7]{3.1}). This condition has two parts, a stratification and a local model that replaces local constancy, and the complexes we use must in addition have finite Tor amplitude. As we explain below, each of the three conditions answers a specific need. For \(\mathcal O_E\) we then show that constructibility can be tested on the reduction modulo \(\pi_E\) (Proposition \hyperref[bp34]{3.11}). For a morphism \(g:Z\to X\) and a \(\Lambda_X\)-module \(M\) we write \(M|_Z=g^*M\), a \(\Lambda_Z\)-module.

A stratification cannot be avoided, already for sheaves coming from geometry: for \(j:\mathbf G_m\hookrightarrow\mathbf A^1\), the sheaf \(j_!\Lambda\) is zero at the origin and \(\Lambda\) elsewhere, so it is not locally constant on \(\mathbf A^1\).

Local constancy itself has to be rephrased. For instance, for \(\Lambda=\mathbf Z_\ell\), “locally constant” with value a \(\mathbf Z_\ell\)-module would mean locally a constant sheaf with value the discrete group, but on \(\operatorname{Spec}\mathbf F_q\) such a sheaf has \(H^1=0\) (Section \bpout{1.2}), whereas the answer we want is \(\mathbf Z_\ell\). We therefore ask instead that \(M\) be locally of the form \(\Lambda_U\otimes_{\Lambda(*)}P\), with \(P\) a finitely presented module over the ring \(\Lambda(*)\) of sections over the point, as in the opening of this chapter. This keeps the \(\ell\)-adic topology. For finite \(\Lambda\) it is local constancy, and for \(E\) it says that \(M\) is locally free of finite rank. In this formulation, finite presentation of \(P\) replaces finiteness of the stalks.

Finally, the six operations should commute with change of coefficients, and this needs finite Tor amplitude. Over \(\mathbf Z/\ell^2\) the sheaf \(\mathbf F_\ell\) is constructible, but \(\mathbf F_\ell\otimes^{\mathbf L}_{\mathbf Z/\ell^2}\mathbf F_\ell\) has cohomology \(\mathbf F_\ell\) in every degree \(\leq0\). With the Tor condition imposed, by contrast, the category \(D^b_{\mathrm{ctf}}\) is stable under change of coefficients and under \(\otimes^{\mathbf L}\), and for \(\mathcal O_E\), \(E\) and \(\overline{\mathbf Q}_\ell\) the Tor condition is automatic (Proposition \hyperref[bp40]{3.12}).

% object: definition-constructible-sheaves (blueprint/chap-Constructible-complexes/sec-Constructibility/definition-constructible-sheaves/_section.md)
\phantomsection\label{bp7}%
\begin{bpstatement}{Definition~3.1}{ (Constructible sheaves)}
Let \(\Lambda\) be a commutative condensed ring.

\begin{enumerate}
\item Consider the following condition on a \(\Lambda_X\)-module \(M\): there is a pro-étale cover \(\{U_i\to X\}\) and finitely presented modules \(P_i\) over the ring \(\Lambda(*)\) of sections over the point such that \[
   M|_{U_i}\simeq\Lambda_{U_i}\otimes_{\Lambda(*)}P_i.
   \] Here \(\Lambda(*)\) and \(P_i\) are constant sheaves on \(U_i\), acting through \(\Lambda(*)\to\Lambda_{U_i}\). Equivalently, on each \(U_i\) the module \(M\) has a presentation by a finite matrix with entries in \(\Lambda(*)\).
\item For qcqs \(X\), a \(\Lambda_X\)-module \(M\) is \emph{constructible} if there is a finite stratification of \(X\) by constructible locally closed subschemes such that on each stratum \(S\) the restriction \(M|_S\) satisfies the condition of (1). On an arbitrary scheme this condition is imposed locally on affine opens.
\end{enumerate}

We write \(\Cons(X,\Lambda)\subset\operatorname{Mod}(X_{\mathrm{pro}\acute{e}t},\Lambda)\) for the full subcategory of constructible sheaves, and \(\Dbc(X,\Lambda)\subset
D(X_{\mathrm{pro}\acute{e}t},\Lambda)\) for the full subcategory of bounded complexes with constructible cohomology sheaves. A complex \(K\) has \emph{Tor amplitude in \([a,b]\)} if \(K\otimes^{\mathbf L}_{\Lambda_X}N\) has cohomology only in \([a,b]\) for every sheaf \(N\) of \(\Lambda_X\)-modules. We write \(D^b_{\mathrm{ctf}}(X,\Lambda)\) for the full subcategory of \(\Dbc(X,\Lambda)\) of complexes of finite Tor amplitude.
\end{bpstatement}

% text: prose-03-04
\phantomsection\label{bp8}%
The modules that satisfy the condition of (1) on all of \(X\) are the analogue of local systems, and they are the building blocks on the strata.

% object: definition-lisse-sheaves (blueprint/chap-Constructible-complexes/sec-Constructibility/definition-lisse-sheaves/_section.md)
\phantomsection\label{bp9}%
\begin{bpstatement}{Definition~3.2}{ (Lisse sheaves)}
Let \(\Lambda\) be a commutative condensed ring. A \(\Lambda_X\)-module \(M\) is \emph{lisse} if it satisfies the condition of Definition definition-constructible-sheaves on \(X\) itself, without a stratification. A complex \(K\in D(X_{\mathrm{pro}\acute{e}t},\Lambda)\) is \emph{lisse} if its cohomology sheaves are lisse.
\end{bpstatement}

% text: prose-03-05
\phantomsection\label{bp10}%
For the individual coefficient rings the definitions take a familiar form.

% object: remark-constructible-sheaves-coefficient-rings (blueprint/chap-Constructible-complexes/sec-Constructibility/remark-constructible-sheaves-coefficient-rings/_section.md)
\phantomsection\label{bp11}%
\begin{bpstatement}{Remark~3.3}{ (Constructible sheaves for the individual coefficient rings)}
\begin{enumerate}
\item For Noetherian \(X\), every locally closed subset is constructible, so in Definition definition-constructible-sheaves any finite stratification by locally closed subschemes may be used.
\item For finite \(A\) and qcqs \(X\), \(D^b_{\mathrm{ctf}}(X,A)\) is the pullback of the usual étale category \(D^b_{\mathrm{ctf}}(X_{\acute{e}t},A)\) by Corollary \bpout{2.44}.
\item If \(\Lambda=E\) or \(\overline{\mathbf Q}_\ell\), the ring \(\Lambda(*)\) is a field, so each \(P_i\) is a finite-dimensional vector space, and a lisse module is locally free of finite rank. In other words, a constructible \(E\)- or \(\overline{\mathbf Q}_\ell\)-sheaf is one that is lisse on the strata of a finite stratification [\hyperref[bp93]{BS15}, Definition 6.8.6].
\item If \(\Lambda=\mathcal O_E\), locally on a stratum a constructible \(\mathcal O_{E,X}\)-module is \(P\otimes_{\mathcal O_E}\mathcal O_{E,X}\) for a finitely generated \(\mathcal O_E\)-module \(P\). For example, \(\mathcal O_{E,X}/\pi_E\) is constructible.
\item For the admissible coefficients, the constructible sheaves form an abelian category. For finite \(A\) this is the classical theory of sheaves with finite stalks, and the other cases are Lemma \hyperref[bp56]{3.18}, Corollary \hyperref[bp61]{3.20} and Lemma \bpout{5.19}, which we prove in Subsection \bpout{5.4.2}.
\item For finite \(A\), a lisse \(A_X\)-module is the same as a locally constant sheaf with finite stalks. For \(\Lambda=\mathcal O_E\) a lisse module may have torsion, such as \(\mathcal O_{E,X}/\pi_E\); the lisse \(\mathcal O_{E,X}\)-modules without torsion are the locally free ones of finite rank, since a finitely generated torsion free \(\mathcal O_E\)-module is free.
\end{enumerate}
\end{bpstatement}

% text: prose-03-06
\phantomsection\label{bp12}%
Part (5) is special to the admissible coefficients. For a general condensed ring, kernels of maps between constructible sheaves can fail to be constructible, and then \(\Dbc(X,\Lambda)\) is not triangulated.


% object: sec-Integral-constructible-complexes (blueprint/chap-Constructible-complexes/sec-Integral-constructible-complexes/_section.md), content hidden
\subsection*{\texorpdfstring{3.3\quad Integral constructible complexes}{3.3 Integral constructible complexes}}
\phantomsection\label{bp13}%
\addcontentsline{toc}{subsection}{\texorpdfstring{3.3\quad Integral constructible complexes}{3.3 Integral constructible complexes}}
% text: prose-03-07
\phantomsection\label{bp14}%
Constructibility over \(\mathcal O_E\) is defined by a local condition, but in practice we know a complex \(K\in D(X_{\mathrm{pro}\acute{e}t},\mathcal O_E)\) through its reductions \(K_n=K\otimes^{\mathbf L}_{\mathcal O_{E,X}}(\mathcal O_E/\pi_E^n)_X\) of Definition \bpout{2.61}. The criterion below says that \(K\) is constructible exactly when it is derived complete and \(K_1\) is constructible. Write \(\mathbf F_E=\mathcal O_E/\pi_E\) for the residue field, so that \(K_1\in D(X_{\mathrm{pro}\acute{e}t},\mathbf F_E)\).

To prove the criterion we must produce, on each stratum, one complex of finite free \(\mathcal O_E\)-modules that is locally a model for all the \(K_n\) at once. Each \(K_n\) separately is easy to control: its bounds and its stratification are inherited from \(K_1\) by dévissage along the powers of \(\pi_E\), and on a stratum it is locally isomorphic to a constant perfect complex. The difficulty is that these local isomorphisms are chosen independently at each level, so they have to be made compatible before we pass to the limit. On a weakly contractible cover we do this by a Mittag–Leffler argument, which needs a uniform rate. Such a rate exists because a bounded complex of finite free modules over the discrete valuation ring \(\mathcal O_E\) has finitely generated cohomology, and the torsion of this cohomology is killed by a fixed power of \(\pi_E\). Throughout, we use the étale site only at the finite levels, through Theorem \bpout{2.43}.

% object: example-integral-criterion-point (blueprint/chap-Constructible-complexes/sec-Integral-constructible-complexes/example-integral-criterion-point/_section.md)
\phantomsection\label{bp15}%
\begin{bpstatement}{Example~3.4}{ (The criterion over a point)}
On the spectrum of an algebraically closed field, take for \(K\) the sheaf attached to an \(\mathcal O_E\)-module. The computation is then one with modules, and it shows where the uniform rate comes from. Resolving \(\mathcal O_E/\pi_E^e\) by \([\mathcal O_E\xrightarrow{\pi_E^e}\mathcal O_E]\) in degrees \(-1,0\) and reducing modulo \(\pi_E\), we find

\[
\mathcal O_E\otimes^{\mathbf L}_{\mathcal O_E}\mathbf F_E=\mathbf F_E,
\qquad
\mathcal O_E/\pi_E^e\otimes^{\mathbf L}_{\mathcal O_E}\mathbf F_E
=\bigl[\mathbf F_E\xrightarrow{\,0\,}\mathbf F_E\bigr]
=\mathbf F_E\oplus\mathbf F_E[1]\quad(e\geq1),
\]

since \(\pi_E^e\) becomes zero modulo \(\pi_E\). For \(P=\mathcal O_E^r\oplus\bigoplus_{i=1}^s\mathcal O_E/\pi_E^{e_i}\) the reduction \(K_1\) has \(H^0=\mathbf F_E^{r+s}\) and \(H^{-1}=\mathbf F_E^s\). Thus \(H^{-1}(K_1)\) counts the cyclic torsion summands. For \(K=\bigoplus_{n\geq1}\mathcal O_E/\pi_E^n\) we get \(K_1=\bigoplus_{n\geq1}(\mathbf F_E\oplus\mathbf F_E[1])\). This is not constructible, because its cohomology is an infinite-dimensional vector space in degrees \(0\) and \(-1\). In other words, a constructible \(K_1\) allows only finitely many torsion summands. For \(P\) the single power \(\pi_E^{\max e_i}\) then kills all the torsion, and this is the uniform rate.

The completeness hypothesis cannot be dropped either. The module \(E\) is \(\pi_E\)-divisible, so \(E\otimes^{\mathbf L}\mathcal O_E/\pi_E^n=0\) for all \(n\). Hence \(K_1=0\) is constructible, yet \(E\) is not a finitely generated \(\mathcal O_E\)-module. The criterion does not apply because \(E\) is not derived complete: its derived completion \(R\varprojlim_n(E\otimes^{\mathbf L}\mathcal O_E/\pi_E^n)\) is zero.
\end{bpstatement}

% object: lemma-bounds-finite-reductions (blueprint/chap-Constructible-complexes/sec-Integral-constructible-complexes/lemma-bounds-finite-reductions/_section.md)
\phantomsection\label{bp16}%
\begin{bpstatementit}{Lemma~3.5}{ (Bounds for the finite reductions)}
Let \(X\) be Noetherian and let \(K\in D(X_{\mathrm{pro}\acute{e}t},\mathcal O_E)\) be derived complete with \(K_1\in\Dbc(X,\mathbf F_E)\). Choose \([a,b]\) containing the cohomology of \(K_1\). Then every \(K_n\) lies in \(\Dbc(X,\mathcal O_E/\pi_E^n)\) and has Tor amplitude in \([a,b]\), and \(K\) has cohomology only in \([a,b]\). If the cohomology sheaves of \(K_1\) are locally constant on a stratum, then so are those of every \(K_n\).
\end{bpstatementit}

% proof: proof/lemma-bounds-finite-reductions
\begin{proof}[{Proof}]\phantomsection\label{bp17}%
Everything about \(K_n\) is obtained from \(K_1\) by induction on \(n\), using that \(K_n\) is an extension of \(K_{n-1}\) by \(K_1\). The Tor bound needs an extra argument with minimal resolutions, and the bound on \(K\) itself then follows from the Milnor sequence.

The filtration of \(\mathcal O_E/\pi_E^n\) by the powers of \(\pi_E\) has graded pieces \(\mathbf F_E\). Tensoring its short exact sequences with \(K\) over \(\mathcal O_{E,X}\), we get triangles \(K_1\to K_{n}\to K_{n-1}\), up to the identification \(\pi_E^{n-1}\mathcal O_E/\pi_E^n\simeq\mathbf F_E\). Moreover, complexes with constructible cohomology in \([a,b]\) are closed under extensions in triangles, by the long exact cohomology sequence and the fact that constructible sheaves form a weak Serre subcategory. Induction on \(n\) now shows that \(K_n\) is bounded in \([a,b]\) with constructible cohomology. Since locally constant finite sheaves of modules also form a weak Serre subcategory, the same induction works on a fixed stratum, and a single stratification works for all \(n\). By the comparison theorem for constructible sheaves (Theorem \bpout{2.43}), these finite sheaves and complexes are classical, that is, pulled back from the étale site.

For the Tor bound we use \(K_n\otimes^{\mathbf L}_{\mathcal O_E/\pi_E^n}\mathbf F_E\simeq K_1\). At an étale geometric point \(\bar x\), take a bounded above resolution of the stalk of \(K_n\) by finite free modules over the Artinian local ring \(\mathcal O_E/\pi_E^n\). Cancelling a two-term identity summand whenever a differential matrix has a unit entry, we obtain a minimal resolution, whose differentials have entries in \((\pi_E)\), so they vanish after tensoring with \(\mathbf F_E\). Since stalks commute with derived tensor products, the minimal resolution tensored with \(\mathbf F_E\) computes the stalk of \(K_1\) at \(\bar x\), with zero differential. Its rank in degree \(q\) is therefore \(\dim_{\mathbf F_E}H^q(K_1)_{\bar x}\). This vanishes outside \([a,b]\), and we conclude that the stalk has a bounded finite free resolution in \([a,b]\). Such a bound on stalks gives the same Tor bound on the étale site, and Corollary \bpout{2.44} carries it to the pro-étale site, for Tor against arbitrary sheaves of modules there.

By completeness \(K=R\varprojlim K_n\), so the Milnor sequence (2.2) gives the lower bound and the vanishing above \(b+1\). For \(H^{b+1}(K)\) we use the top degree, where \(H^b(K_n)=H^b(K_{n+1})\otimes_{\mathcal O_E/\pi_E^{n+1}}\mathcal O_E/\pi_E^n\), so the transition maps are epimorphisms. Their \(R^1\varprojlim\) vanishes by the last assertion of Proposition \bpout{2.24}, and this means that \(H^{b+1}(K)=0\).
\end{proof}

% text: prose-03-08
\phantomsection\label{bp18}%
To make the local trivializations of the \(K_n\) compatible, we need two facts about complexes over \(\mathcal O_E\): a compatible system of perfect complexes over the rings \(\mathcal O_E/\pi_E^n\) comes from a single bounded finite free complex, and the images of the cohomology at finite levels stabilize at a uniform rate.

% object: lemma-perfect-complexes-over-oe (blueprint/chap-Constructible-complexes/sec-Integral-constructible-complexes/lemma-perfect-complexes-over-oe/_section.md)
\phantomsection\label{bp19}%
\begin{bpstatementit}{Lemma~3.6}{ (Perfect complexes over $\mathcal O_E$)}
For \(n\ge1\) let \(M_n\in D(\mathcal O_E/\pi_E^n)\) be perfect, with maps \(M_{n+1}\to M_n\) inducing isomorphisms \(M_{n+1}\otimes^{\mathbf L}_{\mathcal O_E/\pi_E^{n+1}}\mathcal O_E/\pi_E^n\simeq M_n\). Then there is a bounded complex \(P^\bullet\) of finite free \(\mathcal O_E\)-modules with isomorphisms \(P^\bullet/\pi_E^n\simeq M_n\) in \(D(\mathcal O_E/\pi_E^n)\), compatible with the transition maps.
\end{bpstatementit}

% proof: proof/lemma-perfect-complexes-over-oe
\begin{proof}[{Proof}]\phantomsection\label{bp20}%
We use that a minimal complex over the local ring \(\mathcal O_E/\pi_E^n\) is determined up to isomorphism of complexes, not only up to homotopy, so that the transition maps can be made into equalities of matrices and passed to the limit.

Represent each \(M_n\) by a minimal finite free complex, obtained by cancelling unit entries in the differential matrices. Its reduction modulo \(\pi_E\) has zero differential, so its rank in degree \(q\) is \(\dim_{\mathbf F_E}H^q(M_1)\), which is independent of \(n\) and nonzero only in finitely many degrees. The given isomorphism \(M_{n+1}/\pi_E^n\simeq M_n\) between minimal complexes is represented by a chain homotopy equivalence \(f\), with homotopy inverse \(g\). Since both complexes have zero differential modulo \(\pi_E\), the relations \(gf-1=dh+hd\) and \(fg-1=dh'+h'd\) show that \(f\) is an isomorphism modulo \(\pi_E\) in each degree. By Nakayama's lemma its matrices are then invertible over \(\mathcal O_E/\pi_E^n\). Lifting them to invertible matrices over \(\mathcal O_E/\pi_E^{n+1}\) and changing bases, we may assume that the differential at level \(n+1\) reduces exactly to the differential at level \(n\), with the given transition in the derived category. Induction on \(n\) gives compatible matrices, and their limits define \(P^\bullet\). Its differential squares to zero because it does so modulo every \(\pi_E^n\), and the identifications with the \(M_n\) are compatible by construction.
\end{proof}

% text: prose-03-09
\phantomsection\label{bp21}%
The second fact is used in the proof of Lemma \hyperref[bp29]{3.9} for the internal Hom complex \((P^\bullet)^\vee\otimes P^\bullet\), whose \(H^0\) at level \(n\) consists of the maps \(P^\bullet/\pi_E^n\to P^\bullet/\pi_E^n\) in the derived category. There it says that a map at level \(n\) which lifts far enough lifts to all levels, and how far is enough depends only on the torsion of the cohomology.

% object: lemma-images-cohomology-finite-levels (blueprint/chap-Constructible-complexes/sec-Integral-constructible-complexes/lemma-images-cohomology-finite-levels/_section.md)
\phantomsection\label{bp22}%
\begin{bpstatementit}{Lemma~3.7}{ (Images of cohomology at finite levels)}
Let \(Q^\bullet\) be a bounded complex of finite free \(\mathcal O_E\)-modules, and let \(\pi_E^e\) kill the torsion of \(H^{q+1}(Q^\bullet)\). Then for \(m\ge n+e\)

\[
\operatorname{im}\bigl(H^q(Q^\bullet/\pi_E^m)\to
H^q(Q^\bullet/\pi_E^n)\bigr)
=\operatorname{im}\bigl(H^q(Q^\bullet)\to
H^q(Q^\bullet/\pi_E^n)\bigr).
\]
\end{bpstatementit}

% text: prose-03-10
\phantomsection\label{bp23}%
Such an \(e\) exists, since \(\mathcal O_E\) is a Noetherian discrete valuation ring and \(H^{q+1}(Q^\bullet)\) is finitely generated.

% proof: proof/lemma-images-cohomology-finite-levels
\begin{proof}[{Proof of Lemma 3.7}]\phantomsection\label{bp24}%
The two-term resolution of \(\mathcal O_E/\pi_E^n\) gives an exact sequence

\[
0\to H^q(Q^\bullet)/\pi_E^n\to H^q(Q^\bullet/\pi_E^n)
\to H^{q+1}(Q^\bullet)[\pi_E^n]\to0.
\]

For the transition from level \(m\) to level \(n\), the induced map on the right-hand terms is multiplication by \(\pi_E^{m-n}\), because the corresponding map of resolutions is the identity in degree zero and multiplication by \(\pi_E^{m-n}\) in degree \(-1\). If \(m-n\ge e\), this map is zero, since \(H^{q+1}(Q^\bullet)[\pi_E^m]\) consists of torsion elements and \(\pi_E^e\) kills them. The image of \(H^q(Q^\bullet/\pi_E^m)\) therefore lies in the left-hand term \(H^q(Q^\bullet)/\pi_E^n\) at level \(n\), which is the image of \(H^q(Q^\bullet)\). The reverse inclusion holds because the map from \(H^q(Q^\bullet)\) factors through level \(m\).
\end{proof}

% text: prose-03-11
\phantomsection\label{bp25}%
The bound \(m\geq n+e\) cannot be improved. To see this, take \(Q^\bullet=[\mathcal O_E\xrightarrow{\pi_E^e}\mathcal O_E]\) in degrees \(0,1\), so that \(H^0(Q^\bullet)=0\) and \(H^1(Q^\bullet)=\mathcal O_E/\pi_E^e\). For \(m\geq e\) the group \(H^0(Q^\bullet/\pi_E^m)\) is \(\pi_E^{m-e}\mathcal O_E/\pi_E^m\), and its image at level \(n\) is generated by \(\pi_E^{m-e}\) modulo \(\pi_E^n\). This image is zero exactly when \(m\geq n+e\), in agreement with the image of \(H^0(Q^\bullet)=0\).

At a single finite level the local structure is a statement about the étale site: a perfect complex with locally constant cohomology is locally constant. On a connected scheme its local model is moreover the same at every point.

% object: lemma-local-constancy-finite-level (blueprint/chap-Constructible-complexes/sec-Integral-constructible-complexes/lemma-local-constancy-finite-level/_section.md)
\phantomsection\label{bp26}%
\begin{bpstatementit}{Lemma~3.8}{ (Locally constant perfect complexes)}
Let \(S\) be a connected Noetherian scheme, let \(A\) be a finite local ring, and let \(L\in D^b(S_{\acute{e}t},A)\) have finite Tor amplitude and cohomology sheaves that are locally constant with finite stalks. Let \(\bar s\) be a geometric point of \(S\) and let \(M\) be a bounded complex of finite free \(A\)-modules with \(M\simeq L_{\bar s}\) in \(D(A)\). Then \(L\) is étale locally on \(S\) isomorphic to the constant complex \(\underline M\).
\end{bpstatementit}

% proof: proof/lemma-local-constancy-finite-level
\begin{proof}[{Proof}]\phantomsection\label{bp27}%
We first show that a model of a stalk spreads out to a neighbourhood, and then use connectedness to see that the model is the same everywhere.

Every stalk \(L_{\bar t}\) has finite Tor amplitude and finite cohomology, hence is perfect over the local ring \(A\), that is, represented by a bounded complex \(M_{\bar t}\) of finite free \(A\)-modules. An isomorphism \(M_{\bar t}\simeq L_{\bar t}\) is then a section at \(\bar t\) of the zeroth cohomology sheaf of \(R\cHom(\underline{M_{\bar t}},L)\); here the stalks of this sheaf are computed stalkwise, since \(\underline{M_{\bar t}}\) is a bounded complex of finite free modules. We extend the section to an étale neighbourhood of \(\bar t\) and lift it there to a morphism \(\underline{M_{\bar t}}\to L\). Its cone has bounded locally constant finite cohomology and zero stalk at \(\bar t\). Hence its finitely many cohomology sheaves vanish after shrinking the neighbourhood, and \(L\) is isomorphic to \(\underline{M_{\bar t}}\) there. In particular, the stalks of \(L\) at all geometric points over the image of this neighbourhood are isomorphic to \(M_{\bar t}\) in \(D(A)\).

For each isomorphism class \(c\) in \(D(A)\), let \(S_c\subset S\) be the set of points at whose geometric points the stalk of \(L\) lies in \(c\). By the previous paragraph and openness of étale morphisms, each \(S_c\) is open. Moreover, the sets \(S_c\) are disjoint and cover \(S\). Since \(S\) is connected, \(S=S_c\) for the class \(c\) of \(M\). Thus every point has an étale neighbourhood on which \(L\simeq\underline M\); compare [Stacks, \href{https://stacks.math.columbia.edu/tag/09BI}{09BI}].
\end{proof}

% text: prose-03-12
\phantomsection\label{bp28}%
On a weakly contractible object the local isomorphisms of the finite levels become global, and the only remaining problem is to choose them compatibly. For this we use the uniform rate of Lemma \hyperref[bp22]{3.7}.

% object: lemma-integral-model-weakly-contractible (blueprint/chap-Constructible-complexes/sec-Integral-constructible-complexes/lemma-integral-model-weakly-contractible/_section.md)
\phantomsection\label{bp29}%
\begin{bpstatementit}{Lemma~3.9}{ (Integral models on weakly contractible schemes)}
Let \(W\) be an affine weakly contractible scheme, let \(K\in D(W_{\mathrm{pro}\acute{e}t},\mathcal O_E)\) be derived complete, and let \(P^\bullet\) be a bounded complex of finite free \(\mathcal O_E\)-modules. Suppose that for every \(n\geq1\) there is an isomorphism \((P^\bullet/\pi_E^n)_W\simeq K_n\) in \(D(W_{\mathrm{pro}\acute{e}t},\mathcal O_E/\pi_E^n)\). Then \(K\simeq\mathcal O_{E,W}\otimes_{\mathcal O_E}P^\bullet\).
\end{bpstatementit}

% proof: proof/lemma-integral-model-weakly-contractible
\begin{proof}[{Proof}]\phantomsection\label{bp30}%
The hypothesis gives an isomorphism \((P^\bullet/\pi_E^n)_W\simeq K_n\) at each level \(n\), but these isomorphisms need not be compatible with the reduction maps, so they do not assemble to a map into \(K=R\varprojlim_nK_n\). What we need is a compatible family of maps, that is, an element of the inverse limit of the Hom groups between the reductions. Such an element lifts to a map \(\mathcal O_{E,W}\otimes_{\mathcal O_E}P^\bullet\to K\), because \(K\) is derived complete. To find one that is an isomorphism at level \(1\), we show that these Hom groups satisfy the Mittag–Leffler condition with a rate \(e\) independent of \(n\). Derived Nakayama then shows that the lift is an isomorphism.

Consider the tower of complexes

\[
\cHom_{(\mathcal O_E/\pi_E^n)_W}\bigl((P^\bullet/\pi_E^n)_W,K_n\bigr)
\in D(W_{\mathrm{pro}\acute{e}t},\mathcal O_E/\pi_E^n),
\]

with transition maps induced by the reduction maps \(K_{n+1}\to K_n\), and the tower of groups \(\Hom_{D(W_{\mathrm{pro}\acute{e}t},\mathcal O_E/\pi_E^n)}\bigl((P^\bullet/\pi_E^n)_W,K_n\bigr)\) obtained by applying \(H^0R\Gamma(W,-)\). Since \(P^\bullet\) is finite free, the first tower is the tower of reductions of \((\mathcal O_{E,W}\otimes_{\mathcal O_E}(P^\bullet)^\vee)\otimes^{\mathbf L}_{\mathcal O_{E,W}}K\).

These Hom groups can be infinite, so the Mittag–Leffler condition has to come from a uniform bound. Let \(Q^\bullet=(P^\bullet)^\vee\otimes_{\mathcal O_E}P^\bullet\) and choose \(e\) such that \(\pi_E^e\) kills the torsion of \(H^1(Q^\bullet)\). Let \(N\ge m\ge n+e\) and choose one isomorphism \((P^\bullet/\pi_E^N)_W\simeq K_N\). Its reductions to levels \(m\) and \(n\) identify the three internal Hom complexes at levels \(N\to m\to n\) with the constant complexes \(Q^\bullet/\pi_E^N\to Q^\bullet/\pi_E^m\to Q^\bullet/\pi_E^n\). Their zeroth cohomology sheaves are therefore the constant sheaves on the groups \(H^0(Q^\bullet/\pi_E^N)\to H^0(Q^\bullet/\pi_E^m)\to H^0(Q^\bullet/\pi_E^n)\). By Lemma \hyperref[bp22]{3.7} the images from levels \(N\) and \(m\) at level \(n\) agree, as subsheaves. Since \(W\) is weakly contractible, taking sections over \(W\) is exact, and so it commutes with taking cohomology and images. Hence the images in the Hom group at level \(n\) from levels \(N\) and \(m\) coincide as well. Since \(e\) does not depend on \(n\), \(m\) or \(N\), we obtain the Mittag–Leffler condition: for every \(n\) the images at level \(n\) are constant from level \(n+e\) on.

To lift, note that each Hom group contains an isomorphism, and its reduction to level \(1\) is again an isomorphism. The stable image at level \(1\) is the image from level \(1+e\), so it contains an invertible class. By the Mittag–Leffler property this class lifts to a compatible family \((\varphi_n)_n\) in the inverse limit of the Hom groups. The Milnor sequence gives a surjection

\[
\begin{aligned}
&H^0R\varprojlim_n R\Gamma\bigl(W,\cHom_{(\mathcal O_E/\pi_E^n)_W}
((P^\bullet/\pi_E^n)_W,K_n)\bigr)\\
&\qquad\longrightarrow
\varprojlim_n\Hom_{D(W_{\mathrm{pro}\acute{e}t},\mathcal O_E/\pi_E^n)}
\bigl((P^\bullet/\pi_E^n)_W,K_n\bigr),
\end{aligned}
\]

whose source is \(\Hom_{D(W_{\mathrm{pro}\acute{e}t},\mathcal O_E)}(\mathcal O_{E,W}\otimes_{\mathcal O_E}P^\bullet,K)\). This holds because derived sections commute with inverse limits, \(K\) is complete, and tensoring with the finite free dual of \(P^\bullet\) commutes with inverse limits. The family \((\varphi_n)\) therefore lifts to a morphism \(\mathcal O_{E,W}\otimes_{\mathcal O_E}P^\bullet\to K\) whose first reduction is an isomorphism. Since both objects are complete, derived Nakayama (Lemma \bpout{2.62}) shows that this morphism is an isomorphism. Here completeness of the source follows termwise from \(\mathcal O_{E,W}=\varprojlim(\mathcal O_E/\pi_E^n)_W\). Indeed, the transition maps are epimorphisms, so this limit agrees with the derived limit by Proposition \bpout{2.24}.
\end{proof}

% object: proposition-local-structure-integral-constructible-complexes (blueprint/chap-Constructible-complexes/sec-Integral-constructible-complexes/proposition-local-structure-integral-constructible-complexes/_section.md)
\phantomsection\label{bp31}%
\begin{bpstatementit}{Proposition~3.10}{ (Local structure of integral constructible complexes)}
Let \(X\) be Noetherian and let \(K\in D(X_{\mathrm{pro}\acute{e}t},\mathcal O_E)\) be derived complete with \(K_1\in\Dbc(X,\mathbf F_E)\). Then there is a finite stratification of \(X\) such that on each stratum \(S\), pro-étale locally,

\[
K|_S\simeq \mathcal O_{E,S}\otimes_{\mathcal O_E}P^\bullet
\]

for a bounded complex \(P^\bullet\) of finite free \(\mathcal O_E\)-modules.
\end{bpstatementit}

% proof: proof/proposition-local-structure-integral-constructible-complexes
\begin{proof}[{Proof}]\phantomsection\label{bp32}%
We first trivialize each \(K_n\) locally, separately for each \(n\), and then make the trivializations compatible. To begin, choose a finite stratification on which all cohomology sheaves of \(K_1\) are locally constant, and refine it by connected components; there are finitely many of these, since the strata are Noetherian.

For every stratum \(S\) the restriction \(K|_S\) is again derived complete with constructible first reduction. Indeed, restriction to an open subscheme preserves limits, and so does restriction to a closed subscheme on the pro-étale site (Proposition \bpout{2.58}). Moreover, both restrictions are exact and commute with the two-term cones defining the reductions. From now on we fix a connected stratum \(S\).

\emph{Local trivializations at each finite level.} By Lemma \hyperref[bp16]{3.5}, each \(K_n|_S\) is classical and of finite Tor amplitude, with locally constant finite cohomology, so we may work on \(S_{\acute{e}t}\). Fix a geometric point \(\bar s\). Then the stalk complex at \(\bar s\) has a minimal finite free model \(M_n\) in the common interval \([a,b]\). By Lemma \hyperref[bp26]{3.8} the complex \(K_n|_S\) is étale locally isomorphic to the constant complex \((M_n)_S\).

On stalks at \(\bar s\), the reduction maps \(K_{n+1}\to K_n\) induce maps \(M_{n+1}\to M_n\) inducing isomorphisms \(M_{n+1}\otimes^{\mathbf L}_{\mathcal O_E/\pi_E^{n+1}}\mathcal O_E/\pi_E^n\simeq M_n\), because the étale stalk functor is exact and symmetric monoidal. By Lemma \hyperref[bp19]{3.6} there is a bounded finite free complex \(P^\bullet\) with \(P^\bullet/\pi_E^n\simeq M_n\). Back on \(S_{\mathrm{pro}\acute{e}t}\), we get for each \(n\) local isomorphisms \((P^\bullet/\pi_E^n)_S\simeq K_n|_S\) in \(D(S_{\mathrm{pro}\acute{e}t},\mathcal O_E/\pi_E^n)\), chosen separately for each \(n\).

\emph{The integral isomorphism.} Choose an affine weakly contractible cover \(W\to S\). For each \(n\), the local isomorphism becomes an isomorphism \((P^\bullet/\pi_E^n)_W\simeq K_n|_W\) on all of \(W\), since by weak contractibility a covering by trivializing charts splits. Moreover, the restriction \(K|_W\) is derived complete, and its reductions are the \(K_n|_W\). By Lemma \hyperref[bp29]{3.9} we conclude that \(K|_W\simeq\mathcal O_{E,W}\otimes_{\mathcal O_E}P^\bullet\).
\end{proof}

% text: prose-03-13
\phantomsection\label{bp33}%
The local structure is the hard half of the following criterion.

% object: proposition-integral-constructibility-criterion (blueprint/chap-Constructible-complexes/sec-Integral-constructible-complexes/proposition-integral-constructibility-criterion/_section.md)
\phantomsection\label{bp34}%
\begin{bpstatementit}{Proposition~3.11}{ (The integral constructibility criterion)}
Let \(X\) be Noetherian. The following conditions on \(K\in D(X_{\mathrm{pro}\acute{e}t},\mathcal O_E)\) are equivalent:

\begin{enumerate}
\item \(K\in\Dbc(X,\mathcal O_E)\);
\item \(K\) is derived complete and \(K_1\in\Dbc(X,\mathbf F_E)\);
\item there is a finite stratification of \(X\) such that on each stratum \(S\), pro-étale locally, \(K|_S\simeq\mathcal O_{E,S}\otimes_{\mathcal O_E}P^\bullet\) for a bounded complex \(P^\bullet\) of finite free \(\mathcal O_E\)-modules;
\item \(K\in D^b_{\mathrm{ctf}}(X,\mathcal O_E)\).
\end{enumerate}
\end{bpstatementit}

% text: prose-03-14
\phantomsection\label{bp35}%
In the classical adic setting, this description of the integral constructible category is [\hyperref[bp93]{BS15}, Proposition 6.8.11] and [Stacks, \href{https://stacks.math.columbia.edu/tag/09C7}{09C7}].

% proof: proof/proposition-integral-constructibility-criterion
\begin{proof}[{Proof of Proposition 3.11}]\phantomsection\label{bp36}%
The substantial implication is (2)\(\Rightarrow\)(3), which is Proposition \hyperref[bp31]{3.10}. The others hold because the complexes \(\mathcal O_{E,S}\otimes_{\mathcal O_E}P^\bullet\) are complete, have finite Tor amplitude, and have constructible reductions. These properties survive extension by zero from strata. We prove (1)\(\Rightarrow\)(2)\(\Rightarrow\)(3)\(\Rightarrow\)(1) and (3)\(\Leftrightarrow\)(4).

\emph{(1) implies (2).} We first treat a single constructible sheaf, locally on a stratum, then extend by zero from the strata, and finally pass to complexes.

A finitely generated \(\mathcal O_E\)-module is a finite sum of copies of \(\mathcal O_E\) and of modules \(\mathcal O_E/\pi_E^r\), so its sheaf has a two-term finite free resolution. The coefficient sheaf \(\mathcal O_{E,U_i}\) on a member \(U_i\) of the cover in Definition definition-constructible-sheaves is derived complete. Indeed, it is the limit of its finite quotients by Lemma \bpout{2.51}, and since the transition maps between these quotients are surjective, this limit agrees with the derived limit by Proposition \bpout{2.24}. As completeness is preserved by finite sums and cones, the modules \(\mathcal O_{E,U_i}\otimes_{\mathcal O_E}P_i\) of Definition definition-constructible-sheaves are complete, and their reductions are bounded locally constant finite complexes. We may check these assertions on a pro-étale cover, because restriction to an object is a localization of the topos, and it preserves limits and the cones defining the reductions.

For a constructible sheaf, choose a finite stratification as in the definition. Then for a stratum \(a:S\hookrightarrow X\), factored as \(a=i\circ j\) as in Section \hyperref[bp3]{3.1}, the open immersion \(j\) is quasi-compact because \(X\) is Noetherian. Extension by zero \(a_!=i_*j_!\) preserves completeness, because it is exact, commutes with the cones defining the reductions, and preserves limits. Here, for \(j_!\) the last point is the assertion about limits in Proposition \bpout{2.58}, see [Stacks, \href{https://stacks.math.columbia.edu/tag/09BP}{09BP}], while for \(i_*\) it holds because \(i_*\) is exact and a right adjoint. Using the finite filtration by open and closed pieces, we see that the whole sheaf is complete with constructible finite reductions. Applying this to the finitely many cohomology sheaves of \(K\) and using the bounded Postnikov triangles, we obtain (2).

\emph{(2) implies (3).} This is Proposition \hyperref[bp31]{3.10}.

\emph{(3) implies (1).} Each of the finitely many strata is quasi-compact, so the local isomorphisms of (3) can be taken on a finite covering by quasi-compact pro-étale objects, and their bounds give a common cohomological interval. Since the map from the constant sheaf with value the discrete ring \(\mathcal O_E\) to \(\mathcal O_{E,S}\) is flat [Stacks, \href{https://stacks.math.columbia.edu/tag/099L}{099L}], the cohomology of \(\mathcal O_{E,S}\otimes_{\mathcal O_E}P^\bullet\) is \(\mathcal O_{E,S}\otimes_{\mathcal O_E}H^q(P^\bullet)\). Here \(H^q(P^\bullet)\) is finitely generated over \(\mathcal O_E\), so these sheaves satisfy the condition of Definition definition-constructible-sheaves. This proves (1).

\emph{(3) is equivalent to (4).} The complexes \(\mathcal O_{E,S}\otimes_{\mathcal O_E}P^\bullet\) have finite Tor amplitude, and since the stratification is finite there is a common interval. Since the restrictions to the strata are exact and jointly conservative, they detect vanishing after tensoring with any sheaf of modules, and (3) implies (4). Conversely, (4) implies (1) by definition of \(D^b_{\mathrm{ctf}}\), and hence (3) by the implications already proved.
\end{proof}

% text: prose-03-15
\phantomsection\label{bp37}%
Lemma \hyperref[bp82]{3.25} compares \(D^b_{\mathrm{ctf}}(X,\mathcal O_E)\) with the classical category of constructible \(\mathcal O_E\)-complexes [Stacks, \href{https://stacks.math.columbia.edu/tag/09C1}{09C1}], which is defined with modules over the constant sheaf with value the discrete ring \(\mathcal O_E\).


% object: sec-Finite-Tor-amplitude (blueprint/chap-Constructible-complexes/sec-Finite-Tor-amplitude/_section.md), content hidden
\subsection*{\texorpdfstring{3.4\quad Finite Tor amplitude}{3.4 Finite Tor amplitude}}
\phantomsection\label{bp38}%
\addcontentsline{toc}{subsection}{\texorpdfstring{3.4\quad Finite Tor amplitude}{3.4 Finite Tor amplitude}}
% text: prose-03-16
\phantomsection\label{bp39}%
For the integral and rational rings the finite Tor amplitude condition is automatic.

% object: proposition-ctf-equals-bc (blueprint/chap-Constructible-complexes/sec-Finite-Tor-amplitude/proposition-ctf-equals-bc/_section.md)
\phantomsection\label{bp40}%
\begin{bpstatementit}{Proposition~3.12}{ (Constructible complexes have finite Tor amplitude for $\ell$-adic coefficients)}
For \(\Lambda=\mathbf Z_\ell,\mathcal O_E,\mathbf Q_\ell,E,\overline{\mathbf Q}_\ell\), we have

\[
D^b_{\mathrm{ctf}}(X,\Lambda)=\Dbc(X,\Lambda).
\]
\end{bpstatementit}

% proof: proof/proposition-ctf-equals-bc
\begin{proof}[{Proof}]\phantomsection\label{bp41}%
For \(\mathcal O_E\), and hence for \(\mathbf Z_\ell\), this is the equivalence of (1) and (4) in the criterion above. For the fields we show that every constructible sheaf is flat, which bounds the Tor amplitude of a bounded complex by its cohomological amplitude.

For a field \(\Lambda(*)\), a constructible sheaf \(M\) is, on each stratum and locally on a cover \(\{U_i\}\), of the form \(\Lambda_{U_i}\otimes_{\Lambda(*)}P_i\) with \(P_i\) a finite-dimensional vector space, so it is locally a finite free \(\Lambda_{U_i}\)-module. Restriction to the strata is exact and the restrictions are jointly conservative, as in the proof of the criterion. Hence every constructible sheaf has Tor amplitude in \([0,0]\). Finally, a bounded complex whose finitely many cohomology sheaves have finite Tor amplitude has finite Tor amplitude, by the triangles of its canonical filtration.
\end{proof}

% text: prose-03-17
\phantomsection\label{bp42}%
For a finite ring that is not regular the condition is needed, as the \(\mathbf Z/\ell^2\)-module \(\mathbf F_\ell\) of Section \hyperref[bp5]{3.2} shows. Moreover, finite Tor amplitude is a condition on the whole complex, not on its cohomology sheaves. The category \(D^b_{\mathrm{ctf}}\) is thick, but in general it is not stable under standard truncations. For example, for \(A=\mathbf Z/\ell^2\) the complex \([A\xrightarrow{\ell}A]\) is perfect, whereas its nonzero cohomology modules are \(\mathbf F_\ell\), which are not perfect. For this reason, the compactness proofs for finite \(A\) work with bounded flat constructible models, and standard truncations are taken in \(\Dbc\). The models exist by the next lemma.

% object: lemma-flat-models-generators (blueprint/chap-Constructible-complexes/sec-Finite-Tor-amplitude/lemma-flat-models-generators/_section.md)
\phantomsection\label{bp43}%
\begin{bpstatementit}{Lemma~3.13}{ (Flat models and generators)}
Every \(C\in D^b_{\mathrm{ctf}}(X,A)\) on \(X_{\mathrm{pro}\acute{e}t}\) is represented by a bounded complex of flat constructible \(A\)-sheaves. A flat constructible sheaf has a finite stratification on whose strata it is locally of the form \(\underline M\) with \(M\) a finite projective \(A\)-module. Consequently, \(D^b_{\mathrm{ctf}}(X,A)\) is generated, using shifts, finite extensions and direct summands, by sheaves \(a_!L\), where \(a:S\hookrightarrow X\) is locally closed and \(L\) is locally of the form \(\underline M\) with \(M\) a finite projective \(A\)-module.
\end{bpstatementit}

% proof: proof/lemma-flat-models-generators
\begin{proof}[{Proof}]\phantomsection\label{bp44}%
The flat model exists on the étale site by a classical result, and we pull it back. The generators then come from two filtrations: the stupid filtration of the model, and the filtration of each term by strata.

By Corollary \bpout{2.44}, \(C=\nu_X^*C_{\acute{e}t}\) with \(C_{\acute{e}t}\in D^b_{\mathrm{ctf}}(X_{\acute{e}t},A)\), where \(\nu_X^*:D(X_{\acute{e}t},A)\to D(X_{\mathrm{pro}\acute{e}t},A)\) is the comparison inverse image. By [Stacks, \href{https://stacks.math.columbia.edu/tag/0F4M}{0F4M}] the complex \(C_{\acute{e}t}\) has a bounded flat model, and since \(\nu_X^*\) is exact on sheaves, its pullback is a bounded flat constructible model on \(X_{\mathrm{pro}\acute{e}t}\). We choose a common stratification on which its finitely many terms are locally constant, and their stalks are then finite projective, because a finite flat module over the Noetherian ring \(A\) is finite projective. If \(F\) is a flat constructible sheaf, then the exact sequences for the successive open and closed strata express \(F\) as a finite extension of the sheaves \(a_!L\) attached to its strata. Similarly, the stupid filtration of the bounded flat complex expresses \(C\) as a finite extension of shifts of its \emph{terms}, and shifts and finite extensions preserve finite Tor amplitude. The same argument proves the corresponding statement on \(X_{\acute{e}t}\), which we use in the proof of Proposition \bpout{5.14} before passing to the pro-étale site.
\end{proof}


% object: sec-Field-coefficients (blueprint/chap-Constructible-complexes/sec-Field-coefficients/_section.md), content hidden
\subsection*{\texorpdfstring{3.5\quad Field coefficients}{3.5 Field coefficients}}
\phantomsection\label{bp45}%
\addcontentsline{toc}{subsection}{\texorpdfstring{3.5\quad Field coefficients}{3.5 Field coefficients}}
% text: prose-03-18
\phantomsection\label{bp46}%
Until the end of this section, \(\Lambda=E\) or \(\overline{\mathbf Q}_\ell\), so that \(D^b_{\mathrm{ctf}}=\Dbc\) and we can argue with constructible cohomology sheaves and standard truncations. Recall (after Definition \hyperref[bp9]{3.2}) that a \(\Lambda_X\)-module is constructible exactly when it is lisse, that is locally free of finite rank for the pro-étale topology, on each stratum of a finite stratification of \(X\). By dévissage we reduce statements about \(\Dbc(X,\Lambda)\) to the simplest constructible sheaves, namely lisse sheaves extended by zero from a stratum, the sheaves \(a_!L\) of Section \hyperref[bp3]{3.1}. This dévissage first splits a complex into its cohomology sheaves by truncation, and then filters each constructible sheaf by strata.

For this dévissage we need constructible sheaves to form an abelian category (Lemma \hyperref[bp56]{3.18}). In the proof we reduce a morphism of lisse \(E\)-sheaves to a morphism of integral lattices, and for \(\overline{\mathbf Q}_\ell\) we first descend the morphism to a finite \(E\). To prepare for this, in the next two lemmas we construct the lattices locally and carry out the descent.

% object: definition-lattice (blueprint/chap-Constructible-complexes/sec-Field-coefficients/definition-lattice/_section.md)
\phantomsection\label{bp47}%
\begin{bpstatement}{Definition~3.14}{ (Lattices)}
Let \(E/\mathbf Q_\ell\) be finite and let \(L\) be a lisse \(E_X\)-module. A \emph{lattice} in \(L\) is a lisse \(\mathcal O_{E,X}\)-submodule \(M\subset L\) with \(M[1/\ell]=L\).
\end{bpstatement}

% object: lemma-lattices-lisse-e-sheaves (blueprint/chap-Constructible-complexes/sec-Field-coefficients/lemma-lattices-lisse-e-sheaves/_section.md)
\phantomsection\label{bp48}%
\begin{bpstatementit}{Lemma~3.15}{ (Lattices in lisse $E$-sheaves)}
Let \(E/\mathbf Q_\ell\) be finite and let \(L\) be a lisse \(E_X\)-module. The sheaf on \(X_{\mathrm{pro}\acute{e}t}\) of lattices in \(L\) is represented by a separated surjective étale \(X\)-scheme, with a universal lattice. In particular, \(L\) admits a lattice étale locally on \(X\).
\end{bpstatementit}

% proof: proof/lemma-lattices-lisse-e-sheaves
\begin{proof}[{Proof}]\phantomsection\label{bp49}%
Where \(L\) is free of rank \(r\), a lattice is the same as a lattice in \(E^r\), so the sheaf of lattices is locally the constant sheaf on the discrete set \(\GL_r(E)/\GL_r(\mathcal O_E)\). The representing scheme is obtained by descending the corresponding disjoint unions of copies of the trivializing objects.

The rank of \(L\) is a locally constant function on \(|X|\). Indeed, on a trivializing covering the subsets of constant rank are open, closed and saturated, and they descend by fpqc descent of open subsets. So \(X\) is the disjoint union of the open and closed subsets on which the rank is constant. Since the sheaf of lattices restricts to the sheaf of lattices on each of them, we may assume that \(L\) has constant rank \(r\).

We now construct the representing scheme. On a pro-étale affine \(V\) where \(L\) is free of rank \(r\), the sheaf of lattices is the constant sheaf on the set

\[
\GL_r(E)/\GL_r(\mathcal O_E)
\]

of lattices in \(E^r\), which is discrete because \(\GL_r(\mathcal O_E)\) is an open subgroup of \(\GL_r(E)\). To glue these constant sheaves, recall from Lemma \bpout{2.51} that on an affine object, sections of \(E_X\) are continuous functions to \(E\), with bounded denominators by quasi-compactness. A change of basis on \(V\times_XV\) is therefore a continuous function to \(\GL_r(E)\). Since the stabilizer of a lattice is open, the image of each fixed lattice under this function is locally constant on \(V\times_XV\). Now consider the disjoint union of copies of \(V\times_XV\) indexed by the lattices in \(E^r\). On each component, the preimages of the target components are then open and closed and form a disjoint open cover, so the induced map from this disjoint union to itself is a morphism of schemes. In this way the cocycle of the transition matrices defines a descent datum of schemes for these disjoint unions of copies of \(V\).

For affine \(V=\operatorname{Spec}B\), the disjoint union of copies of \(V\) indexed by a set \(I\) is ind-quasi-affine, since its finite subunions are affine open subschemes and together they cover it. Since descent along fpqc coverings is effective for ind-quasi-affine morphisms [Stacks, \href{https://stacks.math.columbia.edu/tag/0APK}{0APK}], the descent datum above is effective. Taking a finite affine trivializing covering of each affine open of \(X\) and gluing, we therefore obtain a scheme representing the sheaf of lattices. This scheme is separated and étale over \(X\) by [Stacks, \href{https://stacks.math.columbia.edu/tag/02KU}{02KU}, \href{https://stacks.math.columbia.edu/tag/02VN}{02VN}], and it is surjective because every free module has a lattice. By descent of sheaves of modules, this scheme carries a universal lattice.
\end{proof}

% text: prose-03-19
\phantomsection\label{bp50}%
For example, in rank one the lattices in \(E\) are the \(\pi_E^m\mathcal O_E\) for \(m\in\mathbf Z\), so the discrete set of lattices is \(E^\times/\mathcal O_E^\times\simeq\mathbf Z\). Hence for \(L=E_X\) the representing scheme is \(\coprod_{m\in\mathbf Z}X\).

% object: example-lattices-nodal-curve (blueprint/chap-Constructible-complexes/sec-Field-coefficients/example-lattices-nodal-curve/_section.md)
\phantomsection\label{bp51}%
\begin{bpstatement}{Example~3.16}{ (The scheme of lattices on the nodal curve)}
Let \(E=\mathbf Q_\ell\) and let \(L\) be the sheaf on the nodal curve \(X\) from the start of this chapter, obtained from \(E_Y\) on the chain of lines \(Y\to X\) by letting the generator of \(\mathbf Z\) act by \(\ell\). Over \(Y\) the lattices in \(L|_Y=E_Y\) are the constant lattices \(\ell^m\mathbf Z_\ell\), so over \(Y\) the scheme of lattices is \(\coprod_{m\in\mathbf Z}Y\). The generator of \(\mathbf Z\) moves the copy indexed by \(m\) to the copy indexed by \(m+1\) while translating \(Y\). Hence the scheme of lattices over \(X\), which is the quotient by \(\mathbf Z\), is a single copy of \(Y\). It is surjective and étale over \(X\), as the lemma says, but it has no section. Indeed, a section would trivialize the \(\mathbf Z\)-torsor \(Y\to X\), whereas \(Y\) is connected and \(\coprod_{\mathbf Z}X\) is not. This is the absence of a global lattice seen geometrically, and the universal lattice lives on \(Y\), where \(L\) becomes constant.
\end{bpstatement}

% text: prose-03-20
\phantomsection\label{bp52}%
With lattices we pass from \(E\) to \(\mathcal O_E\). For \(\overline{\mathbf Q}_\ell\) we first need to come down to a finite extension \(E\), and on a quasi-compact scheme this only involves finitely many matrix entries.

% object: lemma-descent-lisse-qbar-sheaves-finite (blueprint/chap-Constructible-complexes/sec-Field-coefficients/lemma-descent-lisse-qbar-sheaves-finite/_section.md)
\phantomsection\label{bp53}%
\begin{bpstatementit}{Lemma~3.17}{ (Descent of lisse $\Qbar$-sheaves)}
A finite collection of lisse \(\QbarX\)-modules and morphisms between them is obtained by extension of scalars from lisse \(E_X\)-modules and their morphisms, for one finite extension \(E/\mathbf Q_\ell\).
\end{bpstatementit}

% proof: proof/lemma-descent-lisse-qbar-sheaves-finite
\begin{proof}[{Proof}]\phantomsection\label{bp54}%
The data are described by finitely many matrices with entries in \(\QbarX\) over finitely many quasi-compact objects, and each entry already lies in some \(E_X\). We make this precise.

Since \(X\) is quasi-compact, we may choose a finite covering by quasi-compact pro-étale affines on which all the modules are free; such refinements exist by [Stacks, \href{https://stacks.math.columbia.edu/tag/0988}{0988}]. Their pairwise intersections \(U_i\times_XU_j\) are quasi-compact since \(X\) is quasi-separated, and quasi-separated since they are subschemes of the affine schemes \(U_i\times U_j\). On a quasi-compact quasi-separated object, taking sections commutes with filtered colimits, because coverings admit finite quasi-compact refinements and the sheaf condition involves only finite limits [Stacks, \href{https://stacks.math.columbia.edu/tag/0739}{0739}]. By Lemma \bpout{2.51} we have \(\QbarX=\varinjlim_EE_X\). Hence the finitely many transition matrices, their inverses, and the matrices of the given morphisms have entries in \(E_X\) for a single finite \(E\). Moreover, their cocycle and compatibility identities already hold over \(E_X\), since \(E_X\to\QbarX\) is injective. Descent of sheaves of modules now gives the required lisse modules and morphisms.
\end{proof}

% text: prose-03-21
\phantomsection\label{bp55}%
With these two lemmas we can show that constructible sheaves with field coefficients form an abelian category, and obtain the filtration by extensions by zero of lisse sheaves that the dévissage uses.

% object: lemma-constructible-sheaves (blueprint/chap-Constructible-complexes/sec-Field-coefficients/lemma-constructible-sheaves/_section.md)
\phantomsection\label{bp56}%
\begin{bpstatementit}{Lemma~3.18}{ (Abelian category of constructible sheaves)}
Let \(\Lambda=E\) or \(\overline{\mathbf Q}_\ell\).

\begin{enumerate}
\item \(\Cons(X,\Lambda)\) is closed under kernels, cokernels and extensions in \(\operatorname{Mod}(X_{\mathrm{pro}\acute{e}t},\Lambda)\). In particular it is an abelian category, and its inclusion into \(\operatorname{Mod}(X_{\mathrm{pro}\acute{e}t},\Lambda)\) is exact.
\item \(\Cons(X,\Lambda)\) is essentially small.
\item Every \(F\in\Cons(X,\Lambda)\) has a finite filtration \(0=F_0\subset F_1\subset\cdots\subset F_m=F\) whose subquotients \(F_k/F_{k-1}\) are of the form \(a_!L\), with \(a:S\hookrightarrow X\) a locally closed immersion and \(L\) a lisse \(\Lambda_S\)-module.
\end{enumerate}
\end{bpstatementit}

% text: prose-03-22
\phantomsection\label{bp57}%
If \(X\) is a curve, part (3) amounts to an exact sequence \(0\to j_!L\to F\to i_*i^*F\to0\), where \(j\) is the inclusion of a dense open subset on which \(F\) is lisse, \(L=j^*F\), and \(i\) is the inclusion of the finitely many closed points of the complement.

% proof: proof/lemma-constructible-sheaves
\begin{proof}[{Proof of Lemma 3.18}]\phantomsection\label{bp58}%
Part (1) is the substantial one. All three closure properties are local on a common stratification, so it suffices to treat lisse sheaves on a single stratum, and there the only difficulty is the kernel of a morphism. Parts (2) and (3), by contrast, are formal.

(1) Given finitely many constructible sheaves, we refine their stratifications to a common one. Then all three closure properties can be checked on a single stratum \(S\).

\emph{Extensions.} Let \(0\to F'\to F\to F''\to0\) be exact with \(F'\) and \(F''\) lisse on \(S\). After passing to a cover, \(F'\) and \(F''\) are free. Since \(F\to F''\) is surjective, after a further cover a basis of \(F''\) lifts to sections of \(F\). Hence the sequence splits locally, and \(F\) is locally free of finite rank.

\emph{Kernels and cokernels.} Locally, a morphism of lisse sheaves is a matrix of continuous functions to \(E\), and a priori its kernel need not be locally free. Once we choose lattices, however, the morphism becomes a constant matrix on a suitable pro-étale cover, because its reductions are maps of finite locally constant sheaves. The kernel and cokernel of a constant matrix over \(E\) are constant of finite rank, and since being lisse is a local property for the pro-étale topology, this suffices.

First let \(\Lambda=E\) and let \(f:L\to L'\) be a morphism of lisse sheaves on \(S\). By Lemma \hyperref[bp48]{3.15} we may replace \(S\) by a connected quasi-compact affine étale neighbourhood \(V\) on which \(L\) and \(L'\) contain lattices \(M\) and \(N\). On a finite quasi-compact pro-étale covering trivializing both lattices, the matrix entries of \(f\) have a common denominator, so \(h=\pi_E^mf\) maps \(M\) into \(N\) for some \(m\ge0\).

The reductions \(M_n=M/\pi_E^nM\) and \(N_n=N/\pi_E^nN\) are finite locally constant étale sheaves, and \(M=\varprojlim M_n\), \(N=\varprojlim N_n\) by Lemma \bpout{2.51}, checked where \(M\) and \(N\) are free. Fix a geometric point \(\bar v\) of \(V\). Then the completed fibres \(M_{\bar v}=\varprojlim(M_n)_{\bar v}\) and \(N_{\bar v}=\varprojlim(N_n)_{\bar v}\) are finite free \(\mathcal O_E\)-modules with continuous \(\pi_1(V,\bar v)\)-actions, and the reductions of \(h\) give an equivariant map \(h_{\bar v}\) between them. To make \(h\) constant, note that each \(M_n\) and \(N_n\) becomes constant on a connected finite étale cover of the connected scheme \(V\), and we may choose these covers to form a tower. Their inverse limit \(V'\to V\) is a connected pro-étale covering, so maps between constant sheaves on \(V'\) are constant. On \(V'\) the lattices \(M\) and \(N\) are identified with the constant sheaves attached to \(M_{\bar v}\) and \(N_{\bar v}\), and \(h\) with the constant matrix \(h_{\bar v}\). Hence on \(V'\) the kernel and the cokernel of \(f=\pi_E^{-m}h\) are the constant sheaves attached to the kernel and the cokernel of \(h_{\bar v}\otimes_{\mathcal O_E}E\). In particular they are lisse.

For \(\Lambda=\overline{\mathbf Q}_\ell\), we descend \(L\), \(L'\) and \(f\) to a finite extension \(E/\mathbf Q_\ell\) by Lemma \hyperref[bp53]{3.17}. Since the functor \(\QbarX\otimes_{E_X}-\) is exact, it takes the kernel and cokernel over \(E\) to those of \(f\). Finally, restriction to strata is exact, so the kernel and cokernel of a morphism of constructible sheaves on \(X\) are constructible.

(2) We choose the sites in a fixed universe. Up to isomorphism, a constructible sheaf is then determined by a finite stratification and, on each stratum, descent data of finite rank on a covering, and these data form a set.

(3) By Noetherian induction we refine the stratification of \(F\) so that there is an open stratum \(j:U\hookrightarrow X\) with closed complement \(i:Z\hookrightarrow X\) that is again a union of strata. The exact sequence \(0\to j_!j^*F\to F\to i_*i^*F\to0\) separates the stratum \(U\) from the rest, and here \(j^*F\) is lisse. Applying the same step to \(i^*F\) on \(Z\) and iterating, we obtain a finite filtration of \(F\) whose subquotients are \(a_!(F|_S)\) for the strata \(a:S\hookrightarrow X\), and each \(F|_S\) is lisse.
\end{proof}

% object: example-running-example (blueprint/chap-Constructible-complexes/sec-Field-coefficients/example-running-example/_section.md)
\phantomsection\label{bp59}%
\begin{bpstatement}{Example~3.19}{ (A running example)}
We use the same few schemes throughout:

\[
\{0\}\xrightarrow{\ i\ }\mathbf A^1_k\xleftarrow{\ j\ }\mathbf G_{m,k},
\qquad \mathbf A^1_k\xrightarrow{\ j_\infty\ }\mathbf P^1_k .
\]

The sheaf \(j_!\Lambda_{\mathbf G_m}\) is constructible for the stratification \(\mathbf A^1=\mathbf G_m\sqcup\{0\}\). Its stalk is \(\Lambda\) at a geometric point of \(\mathbf G_m\) and \(0\) at the origin, so it is not lisse on \(\mathbf A^1\), and the sequence of part (3) is the localization sequence

\[
0\to j_!\Lambda_{\mathbf G_m}\to\Lambda_{\mathbf A^1}\to i_*\Lambda\to0 .
\]

We return to this example for compact supports (after Lemma \bpout{6.4}), for \(i^!\) (Example \bpout{6.25}) and for the trace (Section \bpout{9.6}).
\end{bpstatement}

% text: prose-03-23
\phantomsection\label{bp60}%
\textbf{Exercises.} (1) Compute the stalks of \(a_!L\) for a locally closed immersion \(a\) and a lisse \(L\). (2) For \(K=\Lambda_X[1]\), in which degree is the nonzero cohomology sheaf? (3) Show that \(j_!\Lambda_{\mathbf G_m}\) has no nonzero global sections on \(\mathbf A^1_k\).

For complexes we get the following. Since constructible sheaves form an abelian subcategory, complexes with constructible cohomology are stable under cones, and the filtration of part (3) together with the truncation triangles generates them from the sheaves \(a_!L\).

% object: corollary-coefficient-heart-b (blueprint/chap-Constructible-complexes/sec-Field-coefficients/corollary-coefficient-heart-b/_section.md)
\phantomsection\label{bp61}%
\begin{bpstatementit}{Corollary~3.20}{ (Generators of $D^b_c$ with field coefficients)}
Let \(\Lambda=E\) or \(\overline{\mathbf Q}_\ell\). Then \(\Dbc(X,\Lambda)\) is a triangulated subcategory of \(D(X_{\mathrm{pro}\acute{e}t},\Lambda)\), and every object of \(\Dbc(X,\Lambda)\) is obtained by finitely many cones from objects \(a_!L[n]\), where \(a:S\hookrightarrow X\) is a locally closed immersion, \(L\) is a lisse \(\Lambda_S\)-module and \(n\in\mathbf Z\).
\end{bpstatementit}

% proof: proof/corollary-coefficient-heart-b
\begin{proof}[{Proof}]\phantomsection\label{bp62}%
The first assertion is closure under cones, which we check on cohomology sheaves. Let \(K\to K'\to K''\to K[1]\) be a distinguished triangle with \(K,K'\) in \(\Dbc(X,\Lambda)\). By the long exact cohomology sequence, each \(\mathcal H^i(K'')\) is an extension of \(\ker(\mathcal H^{i+1}K\to\mathcal H^{i+1}K')\) by \(\coker(\mathcal H^iK\to\mathcal H^iK')\). Both are constructible by Lemma lemma-constructible-sheaves, and hence so is \(\mathcal H^i(K'')\). Closure under shifts is clear, since shifts only renumber the cohomology sheaves.

For the second assertion, we argue by induction on the number of nonzero cohomology sheaves of \(K\). If \(i\) is the largest degree with \(\mathcal H^i(K)\ne0\), the truncation triangle \(\tau_{<i}K\to K\to\mathcal H^i(K)[-i]\to\) exhibits \(K\) as the cone of a map \(\mathcal H^i(K)[-i-1]\to\tau_{<i}K\), and \(\tau_{<i}K\) has fewer nonzero cohomology sheaves. This reduces the claim to a single shifted constructible sheaf, and for it the filtration from Lemma lemma-constructible-sheaves yields finitely many short exact sequences, that is, distinguished triangles, with terms \(a_!L\).
\end{proof}


% object: sec-Continuity-of-Ext (blueprint/chap-Constructible-complexes/sec-Continuity-of-Ext/_section.md), content hidden
\subsection*{\texorpdfstring{3.6\quad Continuity of Ext}{3.6 Continuity of Ext}}
\phantomsection\label{bp63}%
\addcontentsline{toc}{subsection}{\texorpdfstring{3.6\quad Continuity of Ext}{3.6 Continuity of Ext}}
% text: prose-03-24
\phantomsection\label{bp64}%
A morphism of \(\overline{\mathbf Q}_\ell\)-complexes should be defined over some finite extension \(E\) of \(\mathbf Q_\ell\). Since \(\QbarX=\varinjlim_EE_X\), this amounts to asking that Hom out of a constructible complex commute with a filtered colimit, where a filtered colimit of complexes means the termwise colimit of a filtered diagram of complexes of sheaves. We prove this continuity for complexes with a common lower bound. In this form it is the pro-étale analogue of the classical fact that cohomology of a quasi-compact quasi-separated scheme commutes with filtered colimits. Besides the descent to a finite coefficient field below, continuity is used in two places. The first is the uniform bound for Ext groups, which is proved first for integral lattices, so we need continuity for \(\mathcal O_E\) as well. The second is the compactness argument for \(E\) and \(\overline{\mathbf Q}_\ell\) (Section \bpout{5.2}).

For cohomology and higher direct images the statement is Corollary \bpout{2.29}, which holds for every sheaf of rings. To reach Hom, recall that a constructible sheaf has a finite filtration by sheaves \(a_!L\) with \(L\) lisse on a stratum \(S\) (Lemma lemma-constructible-sheaves). Using sections with support in a closed subscheme, we turn Hom out of \(a_!L\) into cohomology on \(S\) of a complex built from the target by \(Rj_*\) and restrictions, and the comparison then goes through a hypercohomology spectral sequence. Because of the common lower bound, only finitely many of its terms contribute to a given degree, and each of these terms commutes with the colimit.

% object: lemma-sections-support-closed-subscheme (blueprint/chap-Constructible-complexes/sec-Continuity-of-Ext/lemma-sections-support-closed-subscheme/_section.md)
\phantomsection\label{bp65}%
\begin{bpstatementit}{Lemma~3.21}{ (Sections with closed support)}
Let \(\Lambda\) be admissible, let \(i:Z\hookrightarrow X\) be a closed immersion with open complement \(j:V\hookrightarrow X\), and let \(R\underline\Gamma_Z:D(X_{\mathrm{pro}\acute{e}t},\Lambda)\to D(Z_{\mathrm{pro}\acute{e}t},\Lambda)\) be the local cohomology functor (Definition \bpout{2.59}),

\[
R\underline\Gamma_Z K=i^*\fib(K\longrightarrow Rj_*j^*K).
\]

For a sheaf \(M\) of \(\Lambda_X\)-modules on \(X_{\mathrm{pro}\acute{e}t}\), let \(\underline\Gamma_Z(M)=i^*\ker(M\to j_*j^*M)\), the sheaf on \(Z\) of sections of \(M\) with support in \(Z\).

\begin{enumerate}
\item \(R\underline\Gamma_Z\) is the right derived functor of \(\underline\Gamma_Z\). Consequently the adjunction of Lemma \bpout{2.60} holds in the form \(R\!\operatorname{Hom}_X(i_*N,K)=R\!\operatorname{Hom}_Z(N,R\underline\Gamma_ZK)\).
\item Let \(a=i\circ u:S\hookrightarrow X\) with \(u:S\hookrightarrow Z\) an open immersion, and let \(L\) be a finite locally projective \(\Lambda_S\)-module with dual \(L^\vee\). Then naturally in \(K\in D(X_{\mathrm{pro}\acute{e}t},\Lambda)\) \[
   R\!\operatorname{Hom}_X(a_!L,K)
   =R\Gamma\bigl(S,L^\vee\otimes u^*R\underline\Gamma_ZK\bigr).
   \tag{3.1}
   \]
\end{enumerate}
\end{bpstatementit}

% proof: proof/lemma-sections-support-closed-subscheme
\begin{proof}[{Proof}]\phantomsection\label{bp66}%
Lemma \bpout{2.60} gives the adjunction between \(i_*\) and \(R\underline\Gamma_Z\) on derived categories. What is new here is that \(R\underline\Gamma_Z\) is the right derived functor of \(\underline\Gamma_Z\), and this upgrades the adjunction to \(R\!\operatorname{Hom}\) complexes and yields formula (3.1). For this we show two things: \(\underline\Gamma_Z\) is right adjoint to the exact functor \(i_*\) on sheaves, and \(M\to j_*j^*M\) is surjective for injective \(M\). The second means that on injective sheaves the kernel of this map is also its fibre.

We use that \(i^*\) and \(i_*\) are exact, that \(i^*i_*=\mathrm{id}\), and that every sheaf \(M\) sits in the localization sequence \(0\to j_!j^*M\to M\to i_*i^*M\to0\) [Stacks, \href{https://stacks.math.columbia.edu/tag/09AA}{09AA}, \href{https://stacks.math.columbia.edu/tag/09AH}{09AH}, \href{https://stacks.math.columbia.edu/tag/09BM}{09BM}]. The last fact shows that a sheaf whose restriction to \(V\) vanishes is of the form \(i_*\) of a sheaf on \(Z\). In particular \(\ker(M\to j_*j^*M)=i_*\underline\Gamma_Z(M)\).

\emph{Adjunction on sheaves.} The functor \(\underline\Gamma_Z\) is right adjoint to \(i_*:\operatorname{Mod}(Z_{\mathrm{pro}\acute{e}t},\Lambda)\to\operatorname{Mod}(X_{\mathrm{pro}\acute{e}t},\Lambda)\). Indeed, a map \(i_*N\to M\) composes to zero with \(M\to j_*j^*M\), because \(j^*i_*N=0\), so it factors through \(\ker(M\to j_*j^*M)=i_*\underline\Gamma_Z(M)\). Since \(i_*\) is fully faithful, such factorizations are the maps \(N\to\underline\Gamma_Z(M)\). Moreover, as \(i_*\) is exact, its right adjoint \(\underline\Gamma_Z\) preserves K-injective complexes.

\emph{Surjectivity on injectives.} For an injective sheaf \(I\) the unit \(I\to j_*j^*I\) of the adjunction \(j^*\dashv j_*\) on \(\operatorname{Mod}(X_{\mathrm{pro}\acute{e}t},\Lambda)\) is surjective, because it has a right inverse. To construct it, extend the inclusion \(j_!j^*I\to I\) along \(j_!j^*I\hookrightarrow j_*j^*I\) to a map \(s:j_*j^*I\to I\), which is possible because \(I\) is injective. Then the composite \(j_*j^*I\xrightarrow{s}I\to j_*j^*I\) restricts to the identity on \(V\). Since \(j_*\) is fully faithful, an endomorphism of \(j_*j^*I\) is determined by its restriction to \(V\), and so this composite is the identity.

\emph{Comparison with the fibre.} The right derived functor of \(\underline\Gamma_Z\) agrees with \(i^*\fib(K\to Rj_*j^*K)\). To see this, represent \(K\) by a K-injective complex \(I^\bullet\) with injective terms [Stacks, \href{https://stacks.math.columbia.edu/tag/079P}{079P}]. Then \(\underline\Gamma_Z(I^\bullet)\) computes the right derived functor of \(\underline\Gamma_Z\) at \(K\). On the other side, the complex \(j^*I^\bullet\) is K-injective, because the left adjoint \(j_!\) of \(j^*\) is exact. Hence \(j_*j^*I^\bullet\) represents \(Rj_*j^*K\). By the surjectivity just proved, the sequences

\[
0\to i_*\underline\Gamma_Z(I^n)\to I^n\to j_*j^*I^n\to0
\]

are exact termwise, and they define a distinguished triangle \(i_*\underline\Gamma_Z(I^\bullet)\to K\to Rj_*j^*K\to i_*\underline\Gamma_Z(I^\bullet)[1]\). Applying \(i^*\) and using \(i^*i_*=\mathrm{id}\) identifies \(\underline\Gamma_Z(I^\bullet)\) with \(i^*\fib(K\to Rj_*j^*K)\).

\emph{Part (1).} With \(I^\bullet\) as above, the complex \(\underline\Gamma_Z(I^\bullet)\) is K-injective by the adjunction on sheaves, and it represents \(R\underline\Gamma_ZK\) by the comparison with the fibre. The adjunction on complexes therefore gives

\[
R\!\operatorname{Hom}_X(i_*N,K)=\Hom^\bullet(i_*N,I^\bullet)
=\Hom^\bullet(N,\underline\Gamma_Z(I^\bullet)),
\]

whose right-hand side is \(R\!\operatorname{Hom}_Z(N,R\underline\Gamma_ZK)\).

\emph{Formula (3.1).} Since \(a_!=i_*u_!\), part (1), the adjunction \(u_!\dashv u^*\) and the identity \(\cHom(L,-)=L^\vee\otimes-\) for the finite locally projective sheaf \(L\) give \(R\!\operatorname{Hom}_X(i_*u_!L,K)=R\!\operatorname{Hom}_Z(u_!L,R\underline\Gamma_ZK)
=R\Gamma(S,L^\vee\otimes u^*R\underline\Gamma_ZK)\).
\end{proof}

% text: prose-03-25
\phantomsection\label{bp67}%
Formula (3.1) expresses Hom out of \(a_!L\) through cohomology on the stratum \(S\) and the functors \(i^*\), \(u^*\) and \(Rj_*\). Each of these commutes with filtered colimits of uniformly bounded below complexes, and from this we deduce the continuity of Ext.

% object: lemma-continuity-ext-out-constructible-complex (blueprint/chap-Constructible-complexes/sec-Continuity-of-Ext/lemma-continuity-ext-out-constructible-complex/_section.md)
\phantomsection\label{bp68}%
\begin{bpstatementit}{Lemma~3.22}{ (Continuity of Ext)}
Let \((K_\alpha^\bullet)\) be a filtered diagram of complexes of \(\Lambda_X\)-modules that all vanish below a fixed degree, and let \(M\) be one of the following:

\begin{enumerate}
\item an object of \(\Dbc(X,\Lambda)\), where \(\Lambda=E\) or \(\overline{\mathbf Q}_\ell\);
\item a finite extension of sheaves \(a_!L_0\), with \(a:S\hookrightarrow X\) locally closed and \(L_0\) a finite locally free \(\mathcal O_{E,S}\)-module, where \(\Lambda=\mathcal O_E\).
\end{enumerate}

Then for every integer \(n\) the natural map

\[
\varinjlim_\alpha\Hom(M,K_\alpha^\bullet[n])
\longrightarrow
\Hom\bigl(M,(\varinjlim_\alpha K_\alpha^\bullet)[n]\bigr)
\]

is an isomorphism.
\end{bpstatementit}

% text: prose-03-26
\phantomsection\label{bp69}%
Applied to sheaves, which are complexes in degree \(0\), the lemma says that for \(M\in\Cons(X,\Lambda)\), \(\Lambda=E\) or \(\overline{\mathbf Q}_\ell\), the functors \(\Ext^n(M,-)\) commute with filtered colimits of \(\Lambda_X\)-modules.

% proof: proof/lemma-continuity-ext-out-constructible-complex
\begin{proof}[{Proof of Lemma 3.22}]\phantomsection\label{bp70}%
We reduce to a single generator \(a_!L\) and then use (3.1) to turn Hom into cohomology on the stratum \(S\), where Corollary \bpout{2.29} applies. The common lower bound is used at the end, to compare the abutments of two spectral sequences.

In both cases \(M\) is obtained by finitely many extensions and shifts from sheaves \(a_!L\), with \(a:S\hookrightarrow X\) a locally closed immersion and \(L\) a finite locally free \(\Lambda_S\)-module: in case (1) by Corollary \hyperref[bp61]{3.20}, and in case (2) by assumption. Both sides are cohomological functors of \(M\). By the five lemma, the class of objects \(M\) for which the claim holds is therefore closed under extensions and shifts, so it suffices to treat \(M=a_!L\), with \(a=i\circ
u:S\hookrightarrow Z\hookrightarrow X\) as in Lemma \hyperref[bp65]{3.21}. The only property of \(L\) used below is that \(L^\vee\otimes-\) is exact.

By (3.1), we must show that the cohomology of \(R\Gamma(S,L^\vee\otimes u^*R\underline\Gamma_ZK)\) commutes with the colimit, where \(K=\varinjlim_\alpha K_\alpha^\bullet\). To this end we apply the exact functor \(L^\vee\otimes u^*i^*(-):D(X_{\mathrm{pro}\acute{e}t},\Lambda)\to D(S_{\mathrm{pro}\acute{e}t},\Lambda)\) to the triangle \(i_*R\underline\Gamma_ZK\to K\to Rj_*j^*K\to\) and to the corresponding triangles for the \(K_\alpha^\bullet\). By the long exact sequence and the five lemma, it suffices to treat the two terms

\[
R\Gamma\bigl(S,L^\vee\otimes u^*i^*K\bigr)
\quad\text{and}\quad
R\Gamma\bigl(S,L^\vee\otimes u^*i^*Rj_*j^*K\bigr).
\]

For the first term, Corollary \bpout{2.29} on \(S\) applies, because the functor \(L^\vee\otimes u^*i^*(-)\) is exact, commutes with colimits and preserves the lower bound. For the second term, consider the hypercohomology spectral sequence

\[
E_2^{p,r}=H^p\bigl(S,L^\vee\otimes u^*i^*R^rj_*(j^*K)\bigr)
\Longrightarrow
H^{p+r}\bigl(S,L^\vee\otimes u^*i^*Rj_*j^*K\bigr)
\]

which is functorial in \(K\). The sheaves \(R^rj_*(j^*K_\alpha)\) commute with the colimit by the second part of Corollary \bpout{2.29}, and \(H^p(S,-)\) does by its first part. Hence, after passing to the colimit, the map of spectral sequences is an isomorphism on \(E_2\). The abutments agree as well, since by the common lower bound only finitely many \(E_2^{p,r}\) contribute to a given total degree.
\end{proof}


% object: sec-Descent-to-a-finite-extension (blueprint/chap-Constructible-complexes/sec-Descent-to-a-finite-extension/_section.md), content hidden
\subsection*{\texorpdfstring{3.7\quad Descent to a finite extension}{3.7 Descent to a finite extension}}
\phantomsection\label{bp71}%
\addcontentsline{toc}{subsection}{\texorpdfstring{3.7\quad Descent to a finite extension}{3.7 Descent to a finite extension}}
% text: prose-03-27
\phantomsection\label{bp72}%
The lattice arguments below need the discrete valuation ring \(\mathcal O_E\) with its uniformizer \(\pi_E\), but the valuation ring of \(\overline{\mathbf Q}_\ell\) is not of this kind. Since \(\overline{\mathbf Q}_\ell\) is the filtered union of the finite extensions \(E\) of \(\mathbf Q_\ell\), we show instead that every object and every morphism of \(\Dbc(X,\overline{\mathbf Q}_\ell)\) is defined over some finite \(E\).

For \(C\in D(X_{\mathrm{pro}\acute{e}t},E)\) and a finite extension \(E'/E\), the extensions of scalars are the complexes \(C_{E'}=E'_X\otimes_{E_X}C\) and \(C_{\overline{\mathbf Q}_\ell}=\QbarX\otimes_{E_X}C\). These tensor products need not be derived, since \(E'_X\) and \(\QbarX\) are flat over \(E_X\).

The classical counterpart is the fact that a continuous representation of a profinite group on \(\overline{\mathbf Q}_\ell^r\), with the \(\ell\)-adic topology, takes values in \(\GL_r(E)\) for some finite \(E/\mathbf Q_\ell\). Similarly, a constructible complex, or a morphism between two of them, is described by finitely many matrices with entries in \(\QbarX\), and on a quasi-compact scheme finitely many such entries already lie in a single \(E_X\). We use the next lemma whenever a statement for \(\overline{\mathbf Q}_\ell\) is reduced to one for a finite \(E\), as in Lemma lemma-integral-models and in the uniform Ext bound (Lemma \bpout{5.6}).

% object: lemma-descent-finite-coefficient-field (blueprint/chap-Constructible-complexes/sec-Descent-to-a-finite-extension/lemma-descent-finite-coefficient-field/_section.md)
\phantomsection\label{bp73}%
\begin{bpstatementit}{Lemma~3.23}{ (Descent to a finite coefficient field)}
\begin{enumerate}
\item For \(C,D\in\Dbc(X,E)\) and every integer \(r\), extension of scalars induces an isomorphism \[
   \varinjlim_{E'\supset E}
   \Hom_{D(X_{\mathrm{pro}\acute{e}t},E')}(C_{E'},D_{E'}[r])
   \;\xrightarrow{\ \sim\ }\;
   \Hom_{D(X_{\mathrm{pro}\acute{e}t},\overline{\mathbf Q}_\ell)}(C_{\overline{\mathbf Q}_\ell},D_{\overline{\mathbf Q}_\ell}[r]),
   \] the colimit running over finite extensions \(E'\) of \(E\) in \(\overline{\mathbf Q}_\ell\).
\item Every object of \(\Dbc(X,\overline{\mathbf Q}_\ell)\) is isomorphic to \(K_{\overline{\mathbf Q}_\ell}\) for some \(K\in\Dbc(X,E)\) and some finite \(E/\mathbf Q_\ell\).
\end{enumerate}
\end{bpstatementit}

% proof: proof/lemma-descent-finite-coefficient-field
\begin{proof}[{Proof}]\phantomsection\label{bp74}%
Morphisms descend because \(\QbarX\) is the filtered colimit of the \(E'_X\) and Hom out of a constructible complex commutes with this colimit (Lemma \hyperref[bp68]{3.22}). Objects are built from lisse sheaves by finitely many extensions, each classified by a morphism, and so they descend as well.

(1) Extension of scalars \(E'_X\otimes_{E_X}-\) is left adjoint to restriction of scalars \(D(X_{\mathrm{pro}\acute{e}t},E')\to D(X_{\mathrm{pro}\acute{e}t},E)\). Hence \(\Hom_{D(X_{\mathrm{pro}\acute{e}t},E')}(C_{E'},D_{E'}[r])=\Hom_{D(X_{\mathrm{pro}\acute{e}t},E)}(C,D_{E'}[r])\), and similarly for \(\overline{\mathbf Q}_\ell\). Represent \(D\) by a bounded complex \(D^\bullet\). Then the complexes \(D^\bullet\otimes_{E_X}E'_X\) form a filtered system of complexes of \(E_X\)-modules, all vanishing below the lower bound of \(D^\bullet\), and their colimit is \(D^\bullet\otimes_{E_X}\QbarX\), which represents \(D_{\overline{\mathbf Q}_\ell}\). Applying Lemma \hyperref[bp68]{3.22} to \(C\in\Dbc(X,E)\), we get

\[
\varinjlim_{E'}\Hom_{D(X_{\mathrm{pro}\acute{e}t},E)}(C,D_{E'}[r])
\xrightarrow{\ \sim\ }\Hom_{D(X_{\mathrm{pro}\acute{e}t},E)}(C,D_{\overline{\mathbf Q}_\ell}[r]),
\]

which is the claim after the adjunction.

(2) First let \(K\) be a constructible \(\QbarX\)-module, with filtration \(N_1\subset\cdots\subset N_m=K\) as in Lemma lemma-constructible-sheaves. We descend \(N_k\) by induction on \(k\). The subquotient \(N_k/N_{k-1}=a_!L\) descends, because \(L\) does by Lemma \hyperref[bp53]{3.17} and \(a_!\) commutes with extension of scalars. Suppose \(N_{k-1}\) and \(N_k/N_{k-1}\) descend to \(E\). Then the extension \(0\to N_{k-1}\to N_k\to N_k/N_{k-1}\to0\) gives a distinguished triangle, classified by a morphism \(N_k/N_{k-1}\to N_{k-1}[1]\). By (1) this morphism descends after a finite extension of \(E\). The fibre of the descended morphism, that is, its cone shifted by \([-1]\), is a model of \(N_k\).

For a bounded complex \(K\), we argue in the same way by induction on the number of nonzero cohomology sheaves. Let \(i\) be the largest degree with \(\mathcal H^i(K)\neq0\), and consider the truncation triangle \(\tau_{<i}K\to K\to\mathcal H^i(K)[-i]\to\). Once \(\tau_{<i}K\) and \(\mathcal H^i(K)\) have models over a finite \(E\), the connecting morphism \(\mathcal H^i(K)[-i-1]\to\tau_{<i}K\) descends after a finite extension by (1), and the cone of the descended morphism is a model of \(K\). Each step replaces \(E\) by a finite extension, so after finitely many steps \(K\) is defined over a single finite \(E\).
\end{proof}


% object: sec-Lattices-after-stratification (blueprint/chap-Constructible-complexes/sec-Lattices-after-stratification/_section.md), content hidden
\subsection*{\texorpdfstring{3.8\quad Lattices after stratification}{3.8 Lattices after stratification}}
\phantomsection\label{bp75}%
\addcontentsline{toc}{subsection}{\texorpdfstring{3.8\quad Lattices after stratification}{3.8 Lattices after stratification}}
% text: prose-03-28
\phantomsection\label{bp76}%
We pass from statements with coefficients \(\mathcal O_E/\pi_E^n\) to statements with coefficients \(E\) through integral lattices. A lisse \(E_X\)-module may fail to have a lattice on all of \(X\), but it has one on each stratum of a finite stratification (Lemma \hyperref[bp77]{3.24}), and this is enough for dévissage arguments. In particular, from the stratified statement we obtain the integral generators of Lemma lemma-integral-models. Lemma \hyperref[bp77]{3.24} is also the first step in the proofs of the uniform Ext bound and of the degree bound for proper direct image in Chapter \bpout{5}. On a connected normal scheme, by contrast, there is a global lattice, and lisse \(E\)-sheaves are representations of \(\pi_1(X,\bar x)\); we prove this where it is used, in Proposition \bpout{13.10}.

On the nodal curve \(X\) of Example \hyperref[bp51]{3.16}, the sheaf \(L\) has no global lattice, and one stratification step repairs this. On the open stratum \(X\smallsetminus\{\text{node}\}\simeq\mathbf G_m\) the torsor is trivial, so \(L\) is constant there and has a lattice. The node, with its reduced structure, is the spectrum of a field, which is normal, and hence the restriction of \(L\) to the node has a lattice as well.

% object: lemma-lattices-after-finite-stratification (blueprint/chap-Constructible-complexes/sec-Lattices-after-stratification/lemma-lattices-after-finite-stratification/_section.md)
\phantomsection\label{bp77}%
\begin{bpstatementit}{Lemma~3.24}{ (Lattices after a finite stratification)}
Let \(L\) be a lisse \(E_X\)-module. There is a finite stratification of \(X\) by locally closed subschemes \(S\) such that each \(L|_S\) contains a lisse \(\mathcal O_{E,S}\)-submodule \(L_0\) with \(L_0[1/\ell]=L|_S\).
\end{bpstatementit}

% text: prose-03-29
\phantomsection\label{bp78}%
In the proof we first find a lattice after a surjective étale base change and then push it down by the trace. For the first step we use that the lattices in \(E^r\) form a discrete set on which \(\GL_r(E)\) acts continuously, which makes the sheaf of lattices in \(L\) an étale \(X\)-scheme. For the second step the base change must be finite étale, which we arrange by restricting to the strata of a finite stratification.

% proof: proof/lemma-lattices-after-finite-stratification
\begin{proof}[{Proof of Lemma 3.24}]\phantomsection\label{bp79}%
Let \(Y\to X\) be the étale \(X\)-scheme representing lattices in \(L\) (Lemma \hyperref[bp48]{3.15}), with its universal lattice.

\emph{Pushing down by the trace.} Choose finitely many quasi-compact affine opens of this scheme whose images cover \(X\), and let \(p:V\to X\) be their disjoint union. Then \(p\) is quasi-compact, separated, surjective and étale, and the universal lattice gives a lattice \(M_0\subset p^*L\). By [Stacks, \href{https://stacks.math.columbia.edu/tag/07RY}{07RY}, \href{https://stacks.math.columbia.edu/tag/07RZ}{07RZ}], \(X\) has a finite stratification by locally closed subschemes over each of which \(p\) is finite étale and surjective. On such a stratum \(S\), let \(\operatorname{Tr}:p_*p^*L\to L\) be the trace of the finite étale map \(p\) over \(S\), and define \(L_0\subset L|_S\) as the image of the composite

\[
p_*M_0\longrightarrow p_*p^*L\xrightarrow{\operatorname{Tr}}L.
\]

To see that \(L_0\) is a lattice, we split the finite étale cover locally on \(S\) into \(d\) copies of \(S\). Then \(p_*M_0\) is a direct sum of \(d\) lattices \(M_1,\dots,M_d\) in \(L\), and the trace is their sum, so \(L_0=M_1+\cdots+M_d\). If we further trivialize \(L\) pro-étale locally, each \(M_k\) becomes a locally constant function to the discrete set of lattices in \(E^r\), and after a common refinement each \(M_k\) is the constant sheaf attached to a lattice \(N_k\subset E^r\). Then \(L_0\) is the constant sheaf attached to \(N_1+\cdots+N_d\), which is a finite free \(\mathcal O_E\)-lattice spanning \(E^r\). Hence \(L_0\) is lisse and \(L_0[1/\ell]=L|_S\).
\end{proof}


% object: sec-Integral-models (blueprint/chap-Constructible-complexes/sec-Integral-models/_section.md), content hidden
\subsection*{\texorpdfstring{3.9\quad Integral models}{3.9 Integral models}}
\phantomsection\label{bp80}%
\addcontentsline{toc}{subsection}{\texorpdfstring{3.9\quad Integral models}{3.9 Integral models}}
% text: prose-03-30
\phantomsection\label{bp81}%
To deduce statements over \(E\) and \(\overline{\mathbf Q}_\ell\) from statements over \(\mathcal O_E\), we need integral objects that generate \(\Dbc(X,E)\) and \(\Dbc(X,\overline{\mathbf Q}_\ell)\). We also want to know that our integral constructible complexes are those of the Stacks Project, so that results proved there apply, and we prove both in the next lemma. Recall from Proposition \hyperref[bp34]{3.11} in Section \hyperref[bp13]{3.3} that an object lies in \(\Dbc(X,\mathcal O_E)=D^b_{\mathrm{ctf}}(X,\mathcal O_E)\) precisely when it is derived complete and its reduction modulo \(\pi_E\) is a bounded constructible complex pulled back from \(X_{\acute{e}t}\). The Stacks Project works instead with the category \(\Dc(X,\mathcal O_E)\) of [Stacks, \href{https://stacks.math.columbia.edu/tag/09C1}{09C1}], built from modules over the constant sheaf \(\mathcal O_E\) rather than over \(\mathcal O_{E,X}\).

% object: lemma-integral-models (blueprint/chap-Constructible-complexes/sec-Integral-models/lemma-integral-models/_section.md)
\phantomsection\label{bp82}%
\begin{bpstatementit}{Lemma~3.25}{ (Integral models)}
\begin{enumerate}
\item Restriction of scalars along the map from the constant sheaf \(\mathcal O_E\) to \(\mathcal O_{E,X}\) is an equivalence between the derived complete objects of \(D(X_{\mathrm{pro}\acute{e}t},\mathcal O_E)\) and the derived complete complexes of modules over the constant sheaf \(\mathcal O_E\).
\item An object \(K\) of \(D(X_{\mathrm{pro}\acute{e}t},\mathcal O_E)\) lies in \(\Dbc(X,\mathcal O_E)\) if and only if its restriction of scalars lies in \(\Dc(X,\mathcal O_E)\).
\item \(\Dbc(-,\mathcal O_E)\) is preserved by \(f^*:D(Y_{\mathrm{pro}\acute{e}t},\mathcal O_E)\to D(X_{\mathrm{pro}\acute{e}t},\mathcal O_E)\) for every morphism \(f:X\to Y\), by \(\otimes^{\mathbf L}_{\mathcal O_{E,X}}\) on \(D(X_{\mathrm{pro}\acute{e}t},\mathcal O_E)\), and by \(a_!:D(S_{\mathrm{pro}\acute{e}t},\mathcal O_E)\to D(X_{\mathrm{pro}\acute{e}t},\mathcal O_E)\) for every locally closed immersion \(a:S\hookrightarrow X\).
\item Every object of \(\Dbc(X,E)\) is obtained by finitely many cones and shifts from objects \((a_!L_0)[1/\ell]\), where \(a:S\hookrightarrow X\) is a locally closed immersion and \(L_0\) is a lisse \(\mathcal O_{E,S}\)-module that is locally free of finite rank. Every object of \(\Dbc(X,\overline{\mathbf Q}_\ell)\) is obtained in the same way from the objects \(\QbarX\otimes_{\mathcal O_{E,X}}a_!L_0\), for some finite \(E/\mathbf Q_\ell\).
\end{enumerate}
\end{bpstatementit}

% text: prose-03-31
\phantomsection\label{bp83}%
In (4) each \(L_0\) satisfies the condition of Definition definition-constructible-sheaves, so it lies in \(\Cons(S,\mathcal O_E)\). By (3) the objects \(a_!L_0\) then lie in \(\Dbc(X,\mathcal O_E)\). Their reductions modulo \(\pi_E^n\) are bounded constructible complexes with coefficients in \(\mathcal O_E/\pi_E^n\), which are pulled back from \(X_{\acute{e}t}\) (Corollary \bpout{2.44}).

% proof: proof/lemma-integral-models
\begin{proof}[{Proof of Lemma 3.25}]\phantomsection\label{bp84}%
In (1) and (2) we compare two coefficient sheaves related by a flat map that is an isomorphism modulo every \(\pi_E^n\), so they have the same derived complete objects. For (3) we use the two characterizations of integral constructibility, while for (4) we combine the dévissage of Corollary \hyperref[bp61]{3.20} with the lattices of Lemma \hyperref[bp77]{3.24}.

(1) and (2). The map from the constant sheaf \(\mathcal O_E\) to \(\mathcal O_{E,X}\) is flat and induces isomorphisms modulo every \(\pi_E^n\) [Stacks, \href{https://stacks.math.columbia.edu/tag/099L}{099L}]. Hence by [Stacks, \href{https://stacks.math.columbia.edu/tag/099I}{099I}] restriction of scalars is an equivalence on derived complete objects, which is (1). By [Stacks, \href{https://stacks.math.columbia.edu/tag/09C1}{09C1}], an object of \(\Dc(X,\mathcal O_E)\) is a derived complete complex whose reduction modulo \(\pi_E\) is constructible of finite Tor dimension over \(\mathbf F_E\). Since over the field \(\mathbf F_E\) every bounded complex has finite Tor dimension, this description agrees with the description of \(\Dbc(X,\mathcal O_E)\) in Proposition \hyperref[bp34]{3.11}, and together with (1) this gives (2).

(3) By Proposition \hyperref[bp34]{3.11}, an object of \(D(X_{\mathrm{pro}\acute{e}t},\mathcal O_E)\) lies in \(\Dbc(X,\mathcal O_E)\) if and only if it satisfies condition (2) there: it is derived complete with constructible reduction modulo \(\pi_E\). Equivalently, it satisfies condition (3): there is a finite stratification such that on each stratum \(S\), pro-étale locally, \(K|_S\simeq\mathcal O_{E,S}\otimes_{\mathcal O_E}P^\bullet\) for a bounded complex \(P^\bullet\) of finite free \(\mathcal O_E\)-modules; here the implication from (2) to (3) is Proposition \hyperref[bp31]{3.10}. For each of the three operations we therefore check whichever of these two conditions is convenient.

\emph{Inverse image.} By Lemma \bpout{2.54}, \(f^*\) sends finite free \(\mathcal O_{E,Y}\)-modules to finite free \(\mathcal O_{E,X}\)-modules, so it preserves the local description of condition (3).

\emph{Extension by zero.} The functor \(a_!=i_*j_!\) preserves derived completeness. Indeed, \(j_!\) for a quasi-compact open immersion \(j\) preserves limits by [Stacks, \href{https://stacks.math.columbia.edu/tag/09BP}{09BP}], and \(i_*\) is exact and preserves limits by [Stacks, \href{https://stacks.math.columbia.edu/tag/09AA}{09AA}, \href{https://stacks.math.columbia.edu/tag/09BM}{09BM}]. Moreover, both functors commute with reduction modulo \(\pi_E\) and preserve constructibility of the reductions, so \(a_!\) preserves condition (2).

\emph{Tensor product.} Let \(K,M\in\Dbc(X,\mathcal O_E)\), and choose a stratification for \(K\) as above. Derived completeness and constructibility of the reduction modulo \(\pi_E\) are both preserved by cones, so \(\Dbc(X,\mathcal O_E)\) is closed under cones. Moreover, the localization triangles along the strata exhibit \(K\) as a finite iterated extension of the objects \(a_!a^*K\), for the strata \(a:S\hookrightarrow X\), and the functor \(-\otimes^{\mathbf L}M\) preserves these triangles. It therefore suffices to show that \((a_!a^*K)\otimes^{\mathbf L}M\) lies in \(\Dbc(X,\mathcal O_E)\).

For this we use the projection formula for a locally closed immersion: for \(L\in D(S_{\mathrm{pro}\acute{e}t},\mathcal O_E)\) and \(M\in D(X_{\mathrm{pro}\acute{e}t},\mathcal O_E)\) the canonical map

\[
(a_!L)\otimes^{\mathbf L}M\longrightarrow a_!(L\otimes^{\mathbf L}a^*M)
\]

is an isomorphism. Indeed, write \(a=i\circ j\) with \(j\) open and \(i\) closed, so that \(a_!=i_*j_!\). On \(S\) the map restricts to the identity of \(L\otimes^{\mathbf L}a^*M\), while on every complementary stratum both sides restrict to zero. Since restriction to the strata is exact and jointly conservative by [Stacks, \href{https://stacks.math.columbia.edu/tag/09AD}{09AD}], the map is an isomorphism.

Given the parts on inverse image and extension by zero, we are reduced to showing that \(a^*K\otimes^{\mathbf L}a^*M\) lies in \(\Dbc(S,\mathcal O_E)\). Tensoring with \(a^*K\) preserves derived completeness, because \(a^*K\) is pro-étale locally a bounded complex of finite free modules. Moreover, the reduction modulo \(\pi_E\) of the tensor product is the tensor product of the reductions, and this is constructible.

(4) For \(\overline{\mathbf Q}_\ell\), we first descend to a finite \(E\) by Lemma \hyperref[bp73]{3.23}. By Corollary \hyperref[bp61]{3.20}, every object of \(\Dbc(X,E)\) is obtained by finitely many cones and shifts from sheaves \(a_!L\) with \(L\) lisse. We refine the stratum \(S\) by a finite stratification as in Lemma \hyperref[bp77]{3.24}. This replaces \(a_!L\) by finitely many extensions of such sheaves, and on each new stratum \(L\) contains a lattice \(L_0\) with \(L_0[1/\ell]=L\), so that \(a_!L=(a_!L_0)[1/\ell]\). It remains to see that \(L_0\) is integrally constructible. The reduction of \(L_0\) modulo \(\pi_E^n\) is a locally constant sheaf of finitely presented \(\mathcal O_E/\pi_E^n\)-modules, so it comes from the étale site by [Stacks, \href{https://stacks.math.columbia.edu/tag/099Y}{099Y}]. Thus \(L_0\in\Dbc(S,\mathcal O_E)\), and \(a_!L_0\in\Dbc(X,\mathcal O_E)\) by (3).
\end{proof}

% text: prose-03-32
\phantomsection\label{bp85}%
On \(\mathbf A^1\), part (4) is visible in the running example. There the constant sheaf \(E_{\mathbf A^1}\) is the cone of a map \(i_*E[-1]\to j_!E_{\mathbf G_m}\), and up to shift both terms are of the form \((a_!L_0)[1/\ell]\) with \(L_0\) the constant lattice \(\mathcal O_E\) on the stratum.


% object: sec-Inverting-ell (blueprint/chap-Constructible-complexes/sec-Inverting-ell/_section.md), content hidden
\subsection*{\texorpdfstring{3.10\quad Inverting $\ell$}{3.10 Inverting ell}}
\phantomsection\label{bp86}%
\addcontentsline{toc}{subsection}{\texorpdfstring{3.10\quad Inverting $\ell$}{3.10 Inverting ell}}
% text: prose-03-33
\phantomsection\label{bp87}%
Passing from \(\mathcal O_E\) to \(E\) means inverting \(\ell\), and on derived categories this is a Verdier localization.

% object: proposition-inverting-ell-full-derived-categories (blueprint/chap-Constructible-complexes/sec-Inverting-ell/proposition-inverting-ell-full-derived-categories/_section.md)
\phantomsection\label{bp88}%
\begin{bpstatementit}{Proposition~3.26}{ (Inverting $\ell$ on derived categories)}
Extension of scalars \(E_X\otimes_{\mathcal O_{E,X}}-:
D(X_{\mathrm{pro}\acute{e}t},\mathcal O_E)\to D(X_{\mathrm{pro}\acute{e}t},E)\) identifies \(D(X_{\mathrm{pro}\acute{e}t},E)\) with the Verdier quotient of \(D(X_{\mathrm{pro}\acute{e}t},\mathcal O_E)\) by the complexes \(K\) with \(H^q(K)[1/\ell]=0\) for all \(q\).
\end{bpstatementit}

% proof: proof/proposition-inverting-ell-full-derived-categories
\begin{proof}[{Proof}]\phantomsection\label{bp89}%
A functor with a fully faithful right adjoint is the Verdier localization at its kernel, so we show that restriction of scalars is fully faithful and compute the kernel.

Inverting \(\ell\) is a filtered colimit along multiplication by \(\ell\), hence exact. Therefore \(E_X=\mathcal O_{E,X}[1/\ell]\) is flat over \(\mathcal O_{E,X}\), and \(H^q(E_X\otimes K)=H^q(K)[1/\ell]\). Restriction of scalars is right adjoint to extension of scalars, and the counit is an isomorphism, because \(\ell\) already acts invertibly on an \(E_X\)-module. Restriction is therefore fully faithful, and extension of scalars is the Verdier localization at its kernel. By the formula for \(H^q(E_X\otimes K)\), this kernel consists of the complexes \(K\) with \(H^q(K)[1/\ell]=0\) for all \(q\).
\end{proof}

% text: prose-03-34
\phantomsection\label{bp90}%
The kernel contains more than the complexes killed by a power of \(\ell\), for example \(\bigoplus_{n\ge1}\mathcal O_{E,X}/\ell^n\). It does so even among constructible sheaves: on \(\bigsqcup_{n\ge1}\operatorname{Spec}k\), the sheaf equal to \(\mathcal O_E/\ell^n\) on the \(n\)-th point is killed by inverting \(\ell\) but by no power of \(\ell\). When the underlying space of \(X\) is Noetherian, however, this cannot happen in \(\Dbc(X,\mathcal O_E)\), and there inverting \(\ell\) amounts to tensoring the Hom groups with \(E\).

% object: proposition-inverting-ell (blueprint/chap-Constructible-complexes/sec-Inverting-ell/proposition-inverting-ell/_section.md)
\phantomsection\label{bp91}%
\begin{bpstatementit}{Proposition~3.27}{ (Inverting $\ell$)}
Assume that the underlying topological space of \(X\) is Noetherian (for instance, that \(X\) is of finite type over a field). A complex in \(\Dbc(X,\mathcal O_E)\) that is killed by inverting \(\ell\) is killed by a fixed power of \(\ell\), and extension of scalars induces an equivalence

\[
\Dbc(X,\mathcal O_E)[\ell^{-1}]\xrightarrow{\ \sim\ }\Dbc(X,E),
\]

where the source category has the objects of \(\Dbc(X,\mathcal O_E)\) and Hom groups tensored with \(E\).
\end{bpstatementit}

% proof: proof/proposition-inverting-ell
\begin{proof}[{Proof}]\phantomsection\label{bp92}%
Both assertions are [\hyperref[bp93]{BS15}, Proposition 6.8.14(2)]. We prove the first, which is where the Noetherian hypothesis enters. Let \(K\in\Dbc(X,\mathcal O_E)\). By Definition \hyperref[bp7]{3.1}, there is a finite stratification such that the restriction of each of the finitely many cohomology sheaves of \(K\) to each stratum \(S\) is, on quasi-compact pro-étale charts \(U\), of the form \(\mathcal O_{E,U}\otimes_{\mathcal O_E}P\) with \(P\) finitely generated. Suppose that \(K[1/\ell]=0\). Extension of scalars to \(E_X\) is exact, so \(E_U\otimes_{\mathcal O_E}P=0\), and each such \(P\) is torsion. The underlying space of \(X\) is Noetherian, so each stratum is quasi-compact and finitely many charts suffice. With finitely many \(P\), one power of \(\ell\) kills the restrictions of all cohomology sheaves of \(K\) to all strata. By the localization sequences of Proposition \bpout{2.58}, a larger power kills every cohomology sheaf of \(K\). Since \(K\) is bounded, a further power kills \(K\) itself, by induction along its truncation triangles.

Full faithfulness says that for \(K,L\in\Dbc(X,\mathcal O_E)\) every morphism \(E_X\otimes K\to E_X\otimes L\) is \(\ell^{-c}\) times the image of a morphism \(K\to L\) for some \(c\geq0\), and that a morphism \(K\to L\) becomes zero after extension of scalars only if some power of \(\ell\) kills it. For essential surjectivity, recall from Lemma lemma-integral-models that \(\Dbc(X,E)\) is generated by objects coming from \(\Dbc(X,\mathcal O_E)\). By the description of the Hom groups, a connecting morphism between them lifts to an integral one after multiplication by a power of \(\ell\). Hence finite cones of the generators also have integral models.
\end{proof}



\clearpage\phantomsection\hypertarget{bp-outside}{}%
\section*{Outside this export}
\addcontentsline{toc}{section}{Outside this export}
References marked \ensuremath{{}^\ast} are to parts of the blueprint this document leaves out:
\begin{description}
\item[Section 1.2] Weil cohomology, in 1 Introduction (\bpslug{sec-{}\allowbreak{}Weil-{}\allowbreak{}cohomology})
\item[Proposition 2.24] The Milnor sequence, in Ⅰ The six operations (\bpslug{proposition-{}\allowbreak{}milnor-{}\allowbreak{}sequence})
\item[Corollary 2.29] Continuity of cohomology, in Ⅰ The six operations (\bpslug{corollary-{}\allowbreak{}continuity-{}\allowbreak{}cohomology})
\item[Theorem 2.43] Comparison for constructible sheaves, in Ⅰ The six operations (\bpslug{theorem-{}\allowbreak{}comparison-{}\allowbreak{}constructible-{}\allowbreak{}sheaves})
\item[Corollary 2.44] Finite Tor amplitude under comparison, in Ⅰ The six operations (\bpslug{corollary-{}\allowbreak{}finite-{}\allowbreak{}tor-{}\allowbreak{}amplitude-{}\allowbreak{}under-{}\allowbreak{}comparison})
\item[Lemma 2.51] Comparison with completed coefficients, in Ⅰ The six operations (\bpslug{lemma-{}\allowbreak{}comparison-{}\allowbreak{}completed-{}\allowbreak{}coefficients})
\item[Lemma 2.54] Inverse image of the coefficient sheaves, in Ⅰ The six operations (\bpslug{lemma-{}\allowbreak{}inverse-{}\allowbreak{}image-{}\allowbreak{}coefficient-{}\allowbreak{}sheaves})
\item[Proposition 2.58] Extension by zero and restriction to a closed subscheme, in Ⅰ The six operations (\bpslug{proposition-{}\allowbreak{}extension-{}\allowbreak{}zero-{}\allowbreak{}restriction-{}\allowbreak{}closed-{}\allowbreak{}subscheme})
\item[Definition 2.59] Local cohomology, in Ⅰ The six operations (\bpslug{definition-{}\allowbreak{}local-{}\allowbreak{}cohomology})
\item[Lemma 2.60] Local cohomology is right adjoint to $i_*$, in Ⅰ The six operations (\bpslug{lemma-{}\allowbreak{}local-{}\allowbreak{}cohomology-{}\allowbreak{}right-{}\allowbreak{}adjoint})
\item[Definition 2.61] Derived adic completeness, in Ⅰ The six operations (\bpslug{definition-{}\allowbreak{}derived-{}\allowbreak{}adic-{}\allowbreak{}completeness})
\item[Lemma 2.62] Reduction and derived Nakayama, in Ⅰ The six operations (\bpslug{lemma-{}\allowbreak{}reduction-{}\allowbreak{}derived-{}\allowbreak{}nakayama})
\item[Chapter 5] Construction of $f^!$, in Ⅰ The six operations (\bpslug{chap-{}\allowbreak{}Construction-{}\allowbreak{}of-{}\allowbreak{}f})
\item[Section 5.2] Compactness of constructible complexes, in Ⅰ The six operations (\bpslug{sec-{}\allowbreak{}Compactness-{}\allowbreak{}of-{}\allowbreak{}constructible-{}\allowbreak{}complexes})
\item[Lemma 5.6] Uniform Ext bound, in Ⅰ The six operations (\bpslug{lemma-{}\allowbreak{}uniform-{}\allowbreak{}ext-{}\allowbreak{}bound})
\item[Proposition 5.14] Compactness for finite coefficients, in Ⅰ The six operations (\bpslug{proposition-{}\allowbreak{}compactness-{}\allowbreak{}finite-{}\allowbreak{}coefficients})
\item[Section 5.4] Direct image with compact support, in Ⅰ The six operations (\bpslug{sec-{}\allowbreak{}Direct-{}\allowbreak{}image-{}\allowbreak{}with-{}\allowbreak{}compact-{}\allowbreak{}support})
\item[Subsection 5.4.2] Integral coefficients, in Ⅰ The six operations (\bpslug{subsec-{}\allowbreak{}Integral-{}\allowbreak{}coefficients})
\item[Lemma 5.19] Structure of integral constructible complexes, in Ⅰ The six operations (\bpslug{lemma-{}\allowbreak{}structure-{}\allowbreak{}integral-{}\allowbreak{}constructible-{}\allowbreak{}complexes})
\item[Lemma 6.4] Cohomology of the projective line and the trace, in Ⅰ The six operations (\bpslug{lemma-{}\allowbreak{}cohomology-{}\allowbreak{}projective-{}\allowbreak{}line-{}\allowbreak{}trace})
\item[Example 6.25] Local cohomology at the origin of the affine line, in Ⅰ The six operations (\bpslug{example-{}\allowbreak{}closed-{}\allowbreak{}support})
\item[Section 9.6] Compatibilities, in Ⅱ Duality and the trace formula (\bpslug{sec-{}\allowbreak{}Compatibilities})
\item[Proposition 13.10] Lattices and representations on a normal scheme, in Ⅲ Weights and the Riemann hypothesis (\bpslug{proposition-{}\allowbreak{}lattices-{}\allowbreak{}representations-{}\allowbreak{}normal-{}\allowbreak{}scheme})
\end{description}
\clearpage\phantomsection%
\section*{References}
\addcontentsline{toc}{section}{References}
\bpbibitem{BBD82}{bp94}{A. Beilinson, J. Bernstein, and P. Deligne. \emph{Faisceaux pervers}, volume 100 of \emph{Astérisque}. 1982.}

\bpbibitem{BS15}{bp93}{B. Bhatt and P. Scholze. The pro-étale topology for schemes. \emph{Astérisque}, 369:99–201, 2015. \href{https://people.mpim-bonn.mpg.de/scholze/proetale.pdf}{Author's manuscript}.}

\bpbibitem{Con07}{bp95}{B. Conrad. Deligne's notes on Nagata compactifications. \emph{Journal of the Ramanujan Mathematical Society}, 22:205–257, 2007. \href{https://math.stanford.edu/\~conrad/papers/nagatafinal.pdf}{Author's manuscript}.}

\bpbibitem{Del74}{bp96}{P. Deligne. La conjecture de Weil. I. \emph{Publ. Math. IHÉS}, 43:273–307, 1974.}

\bpbibitem{Del80}{bp97}{P. Deligne. La conjecture de Weil. II. \emph{Publ. Math. IHÉS}, 52:137–252, 1980.}

\bpbibitem{Den22}{bp98}{C. Deninger. There is no “Weil-”cohomology theory with real coefficients for arithmetic curves. \emph{Ann. Sc. Norm. Super. Pisa Cl. Sci.}, 2022. To appear. arXiv:\href{https://arxiv.org/abs/2204.02714}{2204.02714}.}

\bpbibitem{Deu41}{bp99}{M. Deuring. Die Typen der Multiplikatorenringe elliptischer Funktionenkörper. \emph{Abh. Math. Sem. Univ. Hamburg}, 14:197–272, 1941.}

\bpbibitem{Eke90}{bp100}{T. Ekedahl. On the adic formalism. In \emph{The Grothendieck Festschrift, Vol. II}, volume 87 of \emph{Progress in Mathematics}, pages 197–218. Birkhäuser, 1990.}

\bpbibitem{Has36}{bp101}{H. Hasse. Zur Theorie der abstrakten elliptischen Funktionenkörper. I, II, III. \emph{J. Reine Angew. Math.}, 175:55–62, 69–88, 193–208, 1936.}

\bpbibitem{Hum72}{bp102}{J. E. Humphreys. \emph{Introduction to Lie Algebras and Representation Theory}, volume 9 of \emph{Graduate Texts in Mathematics}. Springer, 1972.}

\bpbibitem{Ill94}{bp103}{L. Illusie. Autour du théorème de monodromie locale. \emph{Astérisque}, 223:9–57, 1994. §4. \href{https://www.numdam.org/item/AST_1994__223__9_0/}{Original article}.}

\bpbibitem{Kat88}{bp104}{N. M. Katz. Travaux de Laumon. \emph{Astérisque}, 161–162:105–132, 1988. Séminaire Bourbaki, exp. 691.}

\bpbibitem{KW01}{bp105}{R. Kiehl and R. Weissauer. \emph{Weil conjectures, perverse sheaves and \(\ell\)-adic Fourier transform}, volume 42 of \emph{Ergebnisse der Mathematik}. Springer, 2001.}

\bpbibitem{Lau87}{bp106}{G. Laumon. Transformation de Fourier, constantes d'équations fonctionnelles et conjecture de Weil. \emph{Publ. Math. IHÉS}, 65:131–210, 1987.}

\bpbibitem{Mil86a}{bp107}{J. S. Milne. \href{https://www.jmilne.org/math/xnotes/AVs.pdf}{Abelian Varieties}. In \emph{Arithmetic Geometry}. 1986. Corrected online version.}

\bpbibitem{Mil86b}{bp108}{J. S. Milne. \href{https://www.jmilne.org/math/xnotes/JVs.pdf}{Jacobian Varieties}. In \emph{Arithmetic Geometry}. 1986. Corrected online version.}

\bpbibitem{Ran39}{bp109}{R. A. Rankin. Contributions to the theory of Ramanujan's function \(\tau(n)\) and similar arithmetical functions. II. \emph{Proc. Cambridge Philos. Soc.}, 35:357–372, 1939.}

\bpbibitem{SGA4}{bp110}{M. Artin, A. Grothendieck, and J.-L. Verdier, editors. \emph{Théorie des topos et cohomologie étale des schémas}, volumes 269, 270, 305 of \emph{Lecture Notes in Mathematics}. Springer, 1972–1973. SGA 4.}

\bpbibitem{SGA4½}{bp111}{P. Deligne. \emph{Cohomologie étale}, volume 569 of \emph{Lecture Notes in Mathematics}. Springer, 1977. SGA 4½; \emph{Théorèmes de finitude}, §3, and \emph{Rapport sur la formule des traces}, §3, and \emph{Sommes trigonométriques}, §1.}

\bpbibitem{SGA5}{bp112}{A. Grothendieck, editor. \emph{Cohomologie \(\ell\)-adique et fonctions \(L\)}, volume 589 of \emph{Lecture Notes in Mathematics}. Springer, 1977. SGA 5.}

\bpbibitem{SGA7II}{bp113}{P. Deligne and N. Katz, editors. \emph{Groupes de monodromie en géométrie algébrique}, volume 340 of \emph{Lecture Notes in Mathematics}. Springer, 1973. SGA 7 II; Exposé XIII, (2.1.7).}

\bpbibitem{Stacks}{bp114}{The Stacks Project Authors. \href{https://stacks.math.columbia.edu}{The Stacks Project}. Individual results are cited by their stable tags in the text.}

\bpbibitem{Var05}{bp115}{Y. Varshavsky. Lefschetz–Verdier trace formula and a generalization of a theorem of Fujiwara. 2005. Version 2. arXiv:\href{https://arxiv.org/abs/math/0505564}{math/0505564}. \href{https://arxiv.org/abs/math/0505564v2}{Author's preprint}.}

\bpbibitem{Wei48a}{bp116}{A. Weil. \emph{Sur les courbes algébriques et les variétés qui s'en déduisent}. Hermann, Paris, 1948.}

\bpbibitem{Wei48b}{bp117}{A. Weil. \emph{Variétés abéliennes et courbes algébriques}. Hermann, Paris, 1948.}

\bpbibitem{Wei49}{bp118}{A. Weil. Numbers of solutions of equations in finite fields. \emph{Bull. Amer. Math. Soc.}, 55:497–508, 1949.}

\end{document}
